% concatenated from Part-1-spectral-stability.tex
\documentclass[11pt,a4paper]{article}
\usepackage{fancyhdr}
\usepackage[british]{babel}

\usepackage{amssymb, amsmath, amsthm, amsfonts, enumerate,bm,nicefrac}
\usepackage{amsmath,setspace,scalefnt}
\usepackage[usenames,dvipsnames,svgnames,table]{xcolor}
\usepackage{graphicx,tikz,caption,subcaption}
\usetikzlibrary{patterns}
\usetikzlibrary{shapes}
\usepackage{fullpage}
\usepackage{hyperref}
%\usepackage[notref,notcite]{showkeys}
\usepackage{enumitem,verbatim}%,itemize
\usepackage{multicol}
\usepackage{float}
\usepackage{cleveref}

%\linespread{1.1} 

\usepackage[margin=2.5cm]{geometry}

\definecolor{gray}{rgb}{0.25, 0.25, 0.25}

%\usepackage{showlabels} %%% show the labels

%\usepackage{yufei}


\usepackage{booktabs}  %% For making three-line Table 

\counterwithin{figure}{section}

% ------   Theorem Styles -------
\newtheorem{theorem}{Theorem}[section]
\newtheorem{lemma}[theorem]{Lemma}
\newtheorem{claim}[theorem]{Claim}
\newtheorem{cor}[theorem]{Corollary}
\newtheorem{construction}[theorem]{Construction}
\newtheorem{conj}[theorem]{Conjecture}
\newtheorem{observation}[theorem]{Observation}
\newtheorem*{observation*}{Observation}
\newtheorem{proposition}[theorem]{Proposition}
\newtheorem{problem}[theorem]{Problem}
\newtheorem*{problem*}{Problem}
\newtheorem{question}[theorem]{Question}
\newtheorem*{question*}{Question}
\newtheorem{innerdefinition}{Definition}
\newenvironment{definition*}
  {\renewcommand\theinnerdefinition{}
   \innerdefinition}
  {\endinnerdefinition}

\theoremstyle{definition}
\newtheorem{defn}[theorem]{Definition}
\newtheorem{example}[theorem]{Example}
\newtheorem{setup}[theorem]{Setup}

\theoremstyle{remark}
\newtheorem*{remark}{Remark}
\newtheorem*{notation}{Notation}
\newcounter{propcounter}
%%%%%%%%%%%%%%%%%%%


\newenvironment{poc}{\begin{proof}[Proof of claim]\renewcommand*{\qedsymbol}{$\blacksquare$}}{\end{proof}}

\newcommand{\C}[1]{{\protect\mathcal{#1}}}
\newcommand{\B}[1]{{\bf #1}}
\newcommand{\I}[1]{{\mathbbm #1}}
\renewcommand{\O}[1]{\overline{#1}}
\newcommand{\M}[1]{\mathrm{#1}}
\newcommand{\ser}[1]{\mathsf{#1}}
\newcommand{\G}[1]{\mathfrak{#1}}
\newcommand{\V}[1]{\mathbold{#1}}
\newcommand{\cl}[1]{#1^{\cap}}
\newcommand{\Matrix}[2]{\left(\begin{array}{#1} #2\end{array}\right)}

\newcommand{\ceil}[1]{\lceil #1\rceil}
\newcommand{\floor}[1]{\lfloor #1\rfloor}
\newcommand{\innprod}[1]{\langle #1\rangle}
\newcommand{\me}{{\mathsf e}}
%\newcommand{\dd}{\,\mathrm{d}}
\usepackage{todonotes} 

\newcommand{\eps}{\varepsilon}
\newcommand{\E}{{\rm E}}
\newcommand{\blue}{\textcolor{blue}}

%project specific ones
\newcommand{\ex}{\mathrm{ex}}
\newcommand{\VC}{\mathrm{VC}}
\newcommand{\cP}{\mathcal{P}}
\newcommand{\cQ}{\mathcal{Q}}
\newcommand{\EE}{\mathbb E}

%\newcommand{\fh}{\mathsf{FH}}
%\newcommand{\tw}{\mathrm{tw}}
\newcommand*{\bfrac}[2]{\genfrac{(}{)}{}{}{#1}{#2}}
\newcommand*{\abs}[1]{\lvert#1\rvert}
\newcommand*{\bigabs}[1]{\bigl|#1\bigr|}
\newcommand*{\Bigabs}[1]{\Bigl|#1\Bigr|}


% Comments
\newcommand{\shengtong}[1]{{\color{purple} shengtong: #1}}
\newcommand{\yongtao}[1]{{\color{red} yongtao: #1}}
\newcommand{\hong}[1]{{\color{cyan} Hong: #1}}

% spex

\newcommand{\spex}{\mathrm{spex}_{\mathrm{edge}}}

\newcommand{\red}{\color{red}}

\title{An edge-spectral Erd\H{o}s--Stone--Simonovits theorem\\ and its stability}

\author{
Yongtao Li\thanks{Yau Mathematical Sciences Center, Tsinghua University, Beijing, China. Email: 
\url{yongtao_li@mail.tsinghua.edu.cn}. Supported by the Shuimu Scholar Program at Tsinghua University.} 
\and
Hong Liu\thanks{Extremal Combinatorics and Probability Group (ECOPRO), Institute for Basic Science (IBS), Daejeon, South Korea. Email: \url{hongliu@ibs.re.kr}. Supported by the Institute for Basic Science (IBS-R029-C4).}
\and
Shengtong Zhang\thanks{Department of Mathematics, Stanford University, CA, USA. Email: \url{stzh1555@stanford.edu}. Supported by the Craig Franklin Fellowship in Mathematics at Stanford University.}
}
\date{\today}

\begin{document}
\maketitle


\begin{abstract}
We study the extremal problem that relates the spectral radius $\lambda (G)$ of an $F$-free graph $G$ with its number of edges. 
Firstly, we prove that for any graph $F$ with chromatic number $\chi (F)=r+1\ge 3$, if $G$ is an $F$-free graph on $m$ edges, then $\lambda^2(G)\le {(1-\frac{1}{r} + o(1))2m}$. 
This provides a unified extension of both the Erd\H{o}s--Stone--Simonovits theorem and its vertex-spectral version due to Nikiforov, and confirms a conjecture proposed by Li, Liu and Feng.

We also establish the corresponding edge-spectral stability, showing that if $G$ is an $F$-free graph on $m$ edges with   $\lambda^2(G)=(1- \frac{1}{r} - o(1))2m$, then $G$ differs from a complete bipartite graph by $o(m)$ edges when $r=2$, and $G$ differs from an $r$-partite Tur\'{a}n graph by $o(m)$ edges when $r\ge 3$. 
This extends the classical Erd\H{o}s--Simonovits stability theorem. 

As an application of our method, we improve a result of Zhai, Lin and Shu by showing that if $\lambda (G)>\sqrt{m}$, then there exist two vertices in $G$ that have at least $\frac{1}{2}\sqrt{m} - O(1)$ common neighbors. This bound is the best possible as witnessed by a random construction. 
\end{abstract}

%\begin{center} \emph{In memory of Prof. Vladimir Nikiforov.} \end{center}

%[HL]: Let's display key words and MS code only when the journal requires them.
%{{\bf Key words:}   Extremal graph theory; Eigenvalues; Spectral radius. }

%{{\bf 2010 Mathematics Subject Classification.}  05C35; 05C50.}


%\newpage 
%\tableofcontents % 插入目录




\section{Introduction}

Extremal graph theory investigates how local structural constraints affect global graph parameters with deep connections to other mathematical disciplines. The field originated with foundational questions about how dense a graph can be while avoiding certain forbidden subgraphs, giving rise to its most classical theme: Tur\'{a}n-type problems. The solution to such problems often involves both precise bounds and the characterization of graphs that achieve them; see  \cite{Bollobas78,FS13,Sim2013,GTAC}. 
The stability results describe how nearly extremal graphs -- those close to the maximum size -- must structurally resemble the extremal maximizer; see \cite{Sim1966}. %In this paper, we study stability phenomena from the spectral point of view.
%The study of graph regularity methods, particularly the Szemer\'{e}di regularity lemma and the graph removal lemma, have provided powerful tools for attacking these problems by decomposing graphs into quasi-random components; see \cite{Fur2015}.  

%Beyond Tur\'{a}n theory, extremal graph theory encompasses several other major directions. Ramsey theory explores unavoidable substructures in edge-colored graphs, guaranteeing the appearance of monochromatic or multicolored configurations.  Extremal graph theory has deep connections to other mathematical disciplines, including probability (e.g., random graphs), algebra (spectral methods), and geometry (incidence problems). Its applications are vast, ranging from theoretical computer science, including network design and coding theory, to the distance problems in discrete geometry and even to additive combinatorics; see, e.g., \cite{GTAC}.   



%\subsection{The classical extremal results and their vertex-spectral extensions} 


A graph $G$ is $F$-free if it does not contain a subgraph isomorphic to $F$. A well-known theorem of Tur\'{a}n \cite{Tur1941} 
states that every $n$-vertex $K_{r+1}$-free graph $G$ satisfies $e(G)\le \left(1-\frac{1}{r} \right) \frac{n^2}{2}$; 
see \cite[p. 294]{Bollobas78} and \cite[p. 269]{AZ2014}. 
Since then there are various extensions and generalizations on Tur\'{a}n's theorem; see, e.g., \cite{BT1981,Bon1983}. 
%This is a cornerstone of extremal graph theory, and it has been studied extensively in the literature.  In particular, numerous proofs of Tur\'{a}n's theorem could be found in the literature; see \cite[pp. 294--301]{Bollobas78} and \cite[pp. 269--273]{AZ2014} for more details.  Furthermore, there are various extensions and generalizations on Tur\'{a}n's theorem; see, e.g., \cite{BT1981,Bon1983}. Tur\'{a}n-type questions investigate the maximum possible number of edges in a graph that does not contain certain subgraphs. 
The \emph{Tur\'{a}n number} $\mathrm{ex}(n,F)$ is defined as the maximum number of edges in an $n$-vertex $F$-free graph. 
A celebrated extension of Tur\'{a}n's theorem attributes to the Erd\H{o}s--Stone--Simonovits theorem \cite{ES46,ES66}, which states that the asymptotics of $\mathrm{ex}(n,F)$ is determined by the chromatic number of $F$. In particular, for any graph $F$ with chromatic number $\chi (F)=r+1$, we have 
$\mathrm{ex}(n,F) = \left(  1-\frac{1}{r} + o(1)\right) \frac{n^2}{2}$.  




In this paper, we study the spectral extensions of Erd\H{o}s--Stone--Simonovits theorem.  Spectral graph theory aims to apply the eigenvalues and ranks of matrices associated with graphs to find the structural information of graphs. 
We refer to \cite{Huang2019,JTYZZ2021} for recent breakthroughs and applications. 
The adjacency matrix of an $n$-vertex graph $G$ is defined as $A_G=[a_{i,j}]_{i,j=1}^n$, where $a_{i,j}=1$ if $ij\in E(G)$, and $a_{i,j}=0$ otherwise. Observe that $A_G$ is a real and symmetric nonnegative matrix.  
Let $\lambda (G)$ be the spectral radius of $G$, which is defined as the maximum absolute value of the eigenvalues of $A_G$. 
Spectral extremal graph theory seeks the connections between graph properties and eigenvalues of the associated matrices. As a result of this topic, many of the classical extremal graph results have been generalized to spectral statements. 
For instance, Wilf \cite{Wil1986} proved a spectral extension of Tur\'{a}n's theorem, showing that if $G$ is an $n$-vertex $K_{r+1}$-free graph, then $\lambda (G)\le \left(1- \frac{1}{r} \right)n$. This was generalized by Nikiforov  \cite{Niki2009cpc} in 2009, who established the following spectral counterpart of the Erd\H{o}s--Stone--Simonovits theorem.
\begin{theorem}[Nikiforov~\cite{Niki2009cpc}] 
 \label{thm-spec-ESS} 
If $\chi (F)=r+1$ and $G$ is an $n$-vertex $F$-free graph, then 
\begin{equation*}  \label{eq-spec-ESS}
  \lambda(G) \le 
\left( 1-\frac{1}{r} + o(1)\right) n. 
\end{equation*}
\end{theorem}


\subsection{The Brualdi--Hoffman--Tur\'{a}n problem}

It is well-known that  for every graph $G$ with $m$ edges,\footnote{When we consider the extremal graph problem on a graph in terms of the number of edges, we usually ignore isolated vertices if there are no confusions.} we have
$\lambda (G)< \sqrt{2m}$. 
This bound follows from e.g.~$\lambda^2(G)< \mathrm{Tr}(A^2_G) = \sum_{v\in V}d(v)=2m$.   
%Improving the factor $2$ in (\ref{trivial-bound}), Edwards and Elphick \cite{EE1983} proved that $\lambda (G)\le \sqrt{(1- \nicefrac{1}{\chi})2m}$, where $\chi $ is the chromatic number of $G$; see \cite[Lemma 6]{Hong1995} for an alternative proof.  
This bound can be improved if $G$ has small clique number. In 2002, Nikiforov \cite{Niki2002} proved that if $G$ is $K_{r+1}$-free with $m$ edges, then 
\begin{equation}
    \label{eq-Niki-2002-cpc}
    \lambda (G)\le \sqrt{\left(1- \frac{1}{r} \right)2m}.
\end{equation}
Moreover, the extremal graphs attaining the equality are complete bipartite graphs for $r=2$, or a regular complete $r$-partite graph for every $r\ge 3$; see \cite{Niki2006-walks,Niki2009jctb}.

In contrast with Wilf's theorem, which connects the spectral radius of $G$ with its number of vertices, 
~(\ref{eq-Niki-2002-cpc}) bounds the spectral radius using the number of edges. Note that the edge-version in~(\ref{eq-Niki-2002-cpc}) extends both Tur\'an theorem and Wilf theorem by invoking the fact $\lambda (G)\ge 2m/n$. 
Particularly, if $G$ is an $m$-edge graph with $\lambda (G)> \sqrt{m}$, then $G$ contains a triangle. 
This result was initially studied by Nosal \cite{Nos1970} in 1970, and we refer to \cite{Niki2021, NZ2021, LLZ2024-book-qua} for recent extensions. 

It is worth to point out that the spectral assumption in terms of the number of edges is more general and versatile: the edge-spectral version not only implies that the vertex-spectral version, but also applies to sparse graphs with any edge density. As such, there have been great interest in studying the maximum spectral radius for $F$-free graphs with a given number of edges. One of the first such results was by Brualdi and Hoffman \cite{BH1985} in 1985, where they proved that any $m$-edge graph $G$ satisfies $\lambda (G)\le k-1$ if $m\le {k \choose 2}$. 
 This result was extended by Stanley \cite{Sta1987} who showed that 
 $\lambda (G)\le \frac{1}{2} 
 \left(\sqrt{8m+1} -1 \right).$
 The bound is optimal for complete graphs. 
 %Furthermore, Hong \cite{Hong1988} proved that if $G$ has no isolated vertices, then \[ \lambda (G)\le \sqrt{2m-n+1},\] where the equality holds if and only if $G$ is a star $K_{1,n-1}$ or the complete graph $K_n$. 
See~\cite{HSF2001,Hong1988,Hong1998,Niki2002,ZC2005} for more related results. 
%Let $K_1\vee_t K_s$ be the graph obtained from the complete graph $K_s$ by adding a new vertex joining $t$ vertices of $K_s$. Brualdi and Hoffman \cite{BH1985} conjectured that if $G$ is a graph of size $m$, where $m={s \choose 2} +t$ for some integers $s,t$ with $0\le t \le s-1$, then $\lambda (G)\le \lambda (K_1\vee_t K_s)$, with equality if and only if $G=K_1\vee_t K_s$.  This conjecture was completely solved by Rowlinson \cite{Row1988}. 
Motivated by the aforementioned results of Brualdi,  Hoffman and Tur\'{a}n, it is interesting to consider the following extremal graph problem. 

\begin{problem*}[Brualdi--Hoffman--Tur\'{a}n type problem] 
Given a graph $F$, 
   what is the maximum spectral radius over all $F$-free graphs on $m$ edges without isolated vertices?  
\end{problem*}

%As an important variant of spectral Tur\'{a}n problems, e.g., \Cref{thm-spec-ESS}, the Brualdi--Hoffman--Tur\'{a}n type problem focuses on establishing sharp upper bounds for the spectral radius of graphs that contain $m$ edges while forbidding some specified subgraph $F$. 
This line of investigation has emerged as an active research area in extremal spectral graph theory during recent years.   
%Before our work, there has been a large number of results studying $\spex(m, F)$ for some  particular classes of graphs $F$. 
For example, we refer to \cite{Niki2002,BN2007jctb,Niki2009jctb,Zhang2024,CSZ2024} for the case when $F$ is a clique; 
\cite{Niki2009,ZLS2021,ZS2022dm,LZS2024} when $F$ is a cycle; 
\cite{LLP2023,YLP2024,CY2024,LZZ2024} when $F$ is a fan; 
see also~\cite{Niki2021,NZ2021b,LLZ2024-book-qua,LLZ2025}. 


%\cite{LLP2023,CY2024} study the friendship graph.  \cite{ZW2024dm,CY2025,LZG2025} study the gem graph. \cite{YLP2024,LZZ2024,GL2024} study the fan graph. \cite{NZ2021,LFP2024-triangular,LFP-count-bowtie} study supersaturation phenomena on the number of triangles and \cite{NZ2021b,LLZ2024-book-qua} on $4$-cycles. 

%\section{Main results}

The goal of this paper is to systematically investigate the Brualdi--Hoffman--Tur\'{a}n type problem for general graphs $F$ and establish the corresponding stability result. 




\subsection{Extension of the Erd\H{o}s--Stone--Simonovits theorem}



Our first result resolves the Brualdi--Hoffman--Tur\'{a}n  problem asymptotically for any graph $F$. 

\begin{theorem} \label{conj-LLF}
%Let $F$ be a graph with chromatic number $\chi (F) = r+1\ge 3$. 
If $\chi (F) = r+1 \ge 3$ and $G$ is an $F$-free graph with $m$ edges, then 
$$\lambda^2 (G)\le \left(1-\frac{1}{r} + o(1)\right)2m.$$ 
\end{theorem} 

The bound in \Cref{conj-LLF} is tight 
when $G$ is chosen as an $r$-partite Tur\'{a}n graph with $m$ edges. 
\Cref{conj-LLF} confirms a conjecture of Li, Liu and Feng \cite[Conjecture 1.31]{LLF2022}. Furthermore, combined with the Rayleigh bound $\lambda(G) \geq \frac{2m}{n}$, \Cref{conj-LLF}  implies both the Erd\H{o}s--Stone--Simonovits theorem and its vertex-spectral version in \Cref{thm-spec-ESS}. 



The condition $\chi (F)\ge 3$ in \Cref{conj-LLF} 
is necessary, since \Cref{conj-LLF} does not hold for all bipartite graphs. 
For example, taking $F=K_{2,t}$,  
we can see that $K_{1,m}$ is $F$-free and $\lambda (K_{1,m})= \sqrt{m}$.  
We plan to explore the bipartite case in our future work. 

In 2006, 
Nikiforov \cite{Niki2006-walks} further generalized (\ref{eq-Niki-2002-cpc}). 
An $\ell$-walk is a sequence of $\ell$ 
vertices $v_1,v_2,\ldots ,v_{\ell}$ of $G$ such that 
$v_i$ is adjacent to $v_{i+1}$ for any $i=1,\ldots ,\ell -1$. 
Let $w_{\ell}(G)$ denote the number of $\ell$-walks in $G$. For example, we have $w_1(G)=n$, $w_2(G)=2m$ and
$w_3(G) = \sum_{v\in V}d_v^2$. 


\begin{theorem}[Nikiforov \cite{Niki2006-walks}] \label{thm-Niki-2006}
    If $r\ge 2$ and $G$ is a $K_{r+1}$-free graph, then for every $\ell \ge 1$, we have
    $$\lambda^\ell (G) \leq \left( 1- \frac{1}{r}\right)\cdot w_\ell(G).$$
\end{theorem}
Note that the case $\ell=1$ above is Wilf's result and the case $\ell=2$ is the bound in~\eqref{eq-Niki-2002-cpc}.
%We see from Rayleigh's formula that $\lambda^2(G) \ge \frac{1}{n} \sum_{v\in V} d_v^2$. Wilf's result gives $\lambda (G)\le (1- \frac{1}{r}) n$.  Interestingly, setting $\ell =3$ in \Cref{thm-Niki-2006}, we can get 
%$\lambda^3(G)\le (1- \frac{1}{r} ) \sum_{v\in V}d_v^2$. 
Very recently, Chao and Yu \cite{Chao-Yu-2024} extended \Cref{thm-Niki-2006} 
using the entropy method, connecting powers of spectral radius to homomorphisms of trees in $K_{r+1}$-free graphs. 

Our framework of proof for~\Cref{conj-LLF} is amenable to prove the following result, which generalizes~\Cref{thm-Niki-2006} to graphs $G$ forbidding any general graph $F$. \Cref{thm:walk} is a unified extension of both \Cref{thm-spec-ESS} and \Cref{conj-LLF}. 

\begin{theorem}
    \label{thm:walk}
    If $\chi (F) = r+1\ge 3$ and $G$ is an $F$-free  graph, then 
    $$\lambda^{\ell} (G) \leq \left(1- \frac{1}{r} + o_{\ell, F}(1)\right) \cdot w_\ell(G).$$
\end{theorem} 
 
 




\subsection{An edge-spectral stability result}





In 1966, 
Erd\H{o}s \cite{Erd1966Sta1,Erd1966Sta2} and Simonovits \cite{Sim1966} proved a stability phenomenon for Tur\'{a}n type problems. Given two graphs $G$ and $H$, the \emph{edit distance} $d(G, H)$ is the minimum number of edges we need to change (add/remove) to transfer one graph to the other. The Erd\H{o}s--Simonovits stability states that for any graph $F$ with $\chi (F)=r+1$, any $n$-vertex $F$-free graph $G$ with  
$\left( 1-\frac{1}{r}-o(1) \right)\frac{n^2}{2}$ edges must be of edit distance $o(n^2)$ from the balanced complete $r$-partite Tur\'{a}n graph $T_{n,r}$. 
%\begin{theorem}[Erd\H{o}s--Simonovits \cite{Erd1966Sta1,Erd1966Sta2,Sim1966}] \label{classical-stability}
%Let $F$ be a graph with chromatic number $\chi (F)=r+1$.
%For every $\varepsilon >0$, there exist $\delta >0$ and $n_0\in \mathbb{N}$ 
%such that if $G$ is an $F$-free graph on $n\ge n_0$ vertices and 
%$ e(G)\ge (1- \frac{1}{r} - \delta ) \frac{n^2}{2}$, 
%then $d(G,T_{n,r})\le \varepsilon n^2$. 
%\end{theorem}
We refer the readers to \cite{Fur2015} for a short proof. 

 
The following vertex-spectral stability can be deduced from a result of Nikiforov \cite{Niki2009jgt}. 



\begin{theorem}[Nikiforov \cite{Niki2009jgt}]  
\label{thm-niki-stability}
Let $F$ be a graph with $\chi (F)=r+1\ge 3$. 
For every $\varepsilon >0$, there exist $\delta >0$ 
and $n_0$ such that 
if  $G$ is an $F$-free graph on $n\ge n_0$ vertices  with $\lambda (G) \ge \left(1- \frac{1}{r} -\delta \right)n,$
then $G$ differs from $T_{n,r}$ in at most $\varepsilon n^2$ edges. 
\end{theorem}

\Cref{thm-niki-stability} provides an effective approach for 
studying the spectral problems when forbidding a non-bipartite graph $F$ and has brought several nice applications in recent years.

In this section, we show a similar stability result for the edge-spectral Tur\'{a}n problem. 
For two disjoint vertex sets $A,B$, 
we write $K_{A,B}$ for the complete bipartite graph on the parts $A$ and $B$. 
For a vertex set $C$, we write $T_{C,r}$ 
for an $r$-partite Tur\'{a}n graph on the vertex set $C$. 

\begin{theorem} 
\label{thm:nikiforov-stability}
Let $F$ be a graph with $\chi (F)=r+1\ge 3$.  
For every $\varepsilon > 0 $, 
there exist $\delta>0$ and $m_0$ such that the following holds. If $G$ is an $F$-free graph with $m\ge m_0$ edges and $\lambda^2 (G)\ge (1- \frac{1}{r} - \delta) 2m$, then 
\begin{itemize}
\item[\rm (a)] When $r = 2$, there exist disjoint vertex sets $A, B \subseteq V(G)$ such that $d(G, K_{A, B}) \leq \varepsilon m$.

\item[\rm (b)]
 When $r \geq 3$, there exists a vertex set $C \subseteq V(G)$ and a Tur\'{a}n graph $T_{C, r}$ on $C$ such that $d(G, T_{C, r}) \leq \varepsilon m$. 
 \end{itemize}
\end{theorem}

Similarly, combining with the Rayleigh bound~$\lambda(G)\ge \frac{2m}{n}$, the edge-spectral stability in \Cref{thm:nikiforov-stability} generalizes both the Erd\H{o}s--Simonovits stability and the Nikiforov vertex-spectral stability.
We remark that when $r = 2$, our edge-spectral version in \Cref{thm:nikiforov-stability} differs substantially from the other two stabilities, as the graph in \Cref{thm:nikiforov-stability} can be close to any complete bipartite graph instead of just the balanced complete bipartite graph. Moreover, we emphasize that the disjoint sets $A$ and $B$ are not necessarily a partition of vertices of $G$, e.g., when $G$ contains some isolated vertices or `bad' vertices with sparse edges.

In next subsection, we will provide a quick application of the edge-spectral stability theorem. 
We expect that~\Cref{thm:nikiforov-stability} will serve as a useful tool for solving edge-spectral Tur\'an problems. In subsequent work, we will return to this topic and use \Cref{thm:nikiforov-stability} to derive some interesting consequences for the Brualdi--Hoffman--Tur\'{a}n problem.


\subsection{Application: maximizing common neighbors} 

 

The classical Tur\'{a}n result of F\"{u}redi \cite{Fur1996jcta} states that $\mathrm{ex}(n,K_{2,t+1})=\frac{1}{2}\sqrt{t}n^{3/2} +O(n^{4/3})$. 
In contrast, when forbidding bipartite graphs, the edge-spectral extremal problem has a rather different behavior. 
In 2021, Zhai, Lin and Shu \cite{ZLS2021} proved the following spectral result.  We refer the readers to \cite{Niki2009} for the case $t=1$ and to \cite{ZS2022dm,NZ2021b,LLZ2024-book-qua} for recent results in this direction.

\begin{theorem}[Zhai--Lin--Shu \cite{ZLS2021}]
\label{thm-ZLS-common}
Let $t\ge 2$ and $m\ge 16t^2$ be integers. 
If $G$ is a $K_{2,t+1}$-free graph with $m$ edges, then $\lambda (G)\le \sqrt{m}$, with equality if and only if $G$ is a star.  
\end{theorem}

We point out here that the quadratic dependence on $m$ in terms of $t$ is necessary. Moreover, the above result implies that in any graph $G$ with $\lambda (G)> \sqrt{m}$, we can find a copy of $K_{2,t}$ with $t=\sqrt{m}/4$. On the other hand, there are many examples showing that such a graph $G$ cannot contain a copy of $K_{2,t}$ with $t=\sqrt{m}$. 
Using the edge-spectral stability~\Cref{thm:nikiforov-stability}, we can obtain the following improvement with the optimal bound on $t$.

\begin{theorem} \label{thm-common}
    If $G$ is an $m$-edge graph with 
    $\lambda(G) > \sqrt{m}$, then $G$ contains a copy of $K_{2, t}$ with $t \ge \sqrt{m} / 2 - O(1)$. Moreover, the bound on $t$ is the best possible. 
\end{theorem}






\paragraph{Our approach.} 
To prove \Cref{conj-LLF}, we use the removal lemma to pass from $G$ to a $K_{r+1}$-free subgraph $G'$ by deleting $o(|V(G)|^2)$ edges. If $G$ is dense, we would be done with the help of Nikiforov's result as the removal of these $o(|V(G)|^2)$ edges contributes $o(\sqrt{m})$ to the spectral radius of $G$. Our key novel step in the proof of \Cref{conj-LLF} is to make a reduction to the dense case. To this end,  
we first show that if $G$ is a minimal counterexample, then the weight of every edge $uv$ in $G$ with respect to the Perron--Frobenius vector is large: $x_ux_v=\Omega (m^{-1/2})$ (\Cref{lem-product-entry}). Secondly, 
we bound the $\infty$-norm of the unit Perron--Frobenius eigenvector of $G$ by showing that $\max_{v\in V}\{x_v\}=O(m^{-1/4})$ (\Cref{lem:max}). Combining these two bounds, we can then deduce that the Perron--Frobenius vector of any possible minimal counterexample must be flat, i.e.~each of its coordinates is of order $\Theta(m^{-1/4})$. As a result, this minimal counterexample $G$ must be dense: $|V(G)|=O(\sqrt{m})$. This approach can be extended to prove \Cref{thm:walk}. Different from the edge-spectral setting, we will prove a reduction so that $|V(G)|=O(\lambda (G))$ in the walk-spectral setting for \Cref{thm:walk}.  

For Theorem \ref{thm:nikiforov-stability}, we deliver the proof in three stages. Firstly, we deal with the case when we forbid triangles (\Cref{thm-stability-triangle}). Using the second moment of the spectral radius of $G$, we can remove $o(m)$ edges from $G$ to obtain a bipartite subgraph $G'$ (\Cref{lem-bipartite}). We then show that bipartite graphs with spectral radius close to $\sqrt{m}$ must be close to complete via truncation on its Perron--Frobenius eigenvector (\Cref{lem:close-to-Kab}). Secondly, when we forbid a clique $K_{r+1}$, where $r\ge 3$ (\Cref{thm-stability-clique-4}), we exploit the slackness from the use of Motzkin--Strauss inequality in Nikiforov's bound (\ref{eq-Niki-2002-cpc}). 
This is the bulk of our proof for \Cref{thm:nikiforov-stability}. We analyze the Perron--Frobenius eigenvector of $G$ and use the slackness to show that most of the edges in $G$ has Perron--Frobenius weight around $\Theta(m^{-1/2})$. From this fact, we are able to locate a subset $V'$ of vertices in which every vertex has weight about $\Theta(m^{-1/4})$. Furthermore, majority of the edges lie within $V'$. As a result, the induced subgraph $G'$ on $V'$ satisfies $d(G,G')=o(m)$ and $|V'|\le (1+o(1))\sqrt{2e(G')/(1-\frac{1}{r})}$, which implies $\lambda (G')\ge \left(1- \frac{1}{r} - o(1) \right)|V'|$ (\Cref{lem:nikiforov-stability-clique}). We can then apply the vertex-spectral stability to $G'$ to finish the proof. Finally, when forbidding a general graph $F$ with $\chi (F)=r+1$, we first remove few ``bad'' edges of $G$ that have small contributions to $\lambda (G)$ to obtain a large subgraph $G'$ on vertex set $V'$ such that $x_i'x_j'=\Omega (m^{-1/2})$ for every $ij\in E(G')$ (\Cref{lem:removing-bad-edges}). If $r\ge 3$, we can apply \Cref{lem:max} to show that $\max_{v\in V'}\{x'_v\}=O(m^{-1/4})$, which implies $|V'|=O(\sqrt{m})$.  
Then we can use removal lemma to pass to a $K_{r+1}$-free graph $G''$ by removing $o(|V'|^2)=o(m)$ edges of $G'$. This allows us to reduce the general case to the clique case that we have established previously. For the case when $r=2$, \Cref{lem:max} cannot be applied and 
it is possible that $\max_{v\in V'}\{x_v'\}=\Omega(m^{-1/4})$. 
We need to treat this case separately (\Cref{lem:max-2}). Here we show that $G'$ contains a bipartite subgraph $G''$ with large spectral radius, and we can use \Cref{lem:close-to-Kab} to finish the proof. 

For \Cref{thm-common}, we first show that if 
$\lambda(G)> \sqrt{m}$, then 
$t+ 2m\cdot (\max_{i\in V} x_i)^4 >\frac{1}{2} \sqrt{m}$, where $t$ is the maximum codegree. If $\max_{i\in V} \{x_i\} = O(m^{-1/4})$, then $t\ge \frac{1}{2}\sqrt{m} - O(1)$, as desired. If $\max_{i\in V}\{x_i\} =\Omega (m^{-1/4})$, then we can use the edge-spectral stability (\Cref{lem:max-2}) to get two disjoint vertex sets $A,B\subseteq V(G)$ such that $d(G,K_{A,B})\le \varepsilon m$. Assume that $|A|\le |B|$. As $\lambda(G)> \sqrt{m}$, we can use a structural result (\Cref{lem-Delta}) to show that $|A|\ge 2$. Then by a standard average argument, there exist two vertices of $A$ that have at least $(1-o(1))\sqrt{m}$ common neighbors in $B$, as desired.


\paragraph{Organization.}
In Section \ref{sec:prelim}, we provide some preliminaries. 
In Section \ref{sec:ESS}, 
we give the proofs of Theorems \ref{conj-LLF} and \ref{thm:walk}. In Section \ref{sec:edge-stability}, we present the proof of Theorem \ref{thm:nikiforov-stability}. 
In Section \ref{sec:common}, we give the proof of \Cref{thm-common} and in~\Cref{sec:concluding} we provide some concluding remarks. 
In Appendix \ref{sec-A}, we give an alternative proof of Theorem \ref{thm:nikiforov-stability} in the case when the forbidden graph $F$ is a clique. 


\paragraph{Notation.} 
We write $G=(V,E)$ for a simple and undirected graph with vertex set $V=\{v_1,\ldots ,v_n\}$ and edge set $E=\{e_1,\ldots ,e_m\}$, where we admit 
$n=|V|$ and $m=|E|$. 
We write $uv$ for an edge if $u$ and $v$ are adjacent.  
The degree of a vertex $v\in V(G)$ is denoted by 
$d_v$, and the set of neighbors of $v$ is denoted by $N(v)$. 
The maximum degree of $G$ is denoted by 
$\Delta (G)$. 
Let $w_{\ell} (G)$ be the number of walks in $G$ consisting of $\ell$ vertices (we call it an $\ell$-walk). 
Let $\bm{1}_v$ be the characteristic vector of a vertex $v$.  
For $U\subseteq V(G)$, 
we write $d_U(v)$ for the number of vertices of $U$ that are adjacent to the vertex $v$.  
For two disjoint sets $U,W\subseteq V(G)$, 
we write $e(U)$ for the number of edges with endpoints in $U$, and $e(U,W)$ for the number of edges between $U$ and $W$. 
The edit distance between two graphs $G$ and $H$ is denoted by $d(G,H)$, which is the minimum number of edges we need to change (delete or add) to transform $G$ into $H$.  

The adjacency matrix $A_G$ 
of a graph $G$ is a symmetric matrix of dimension $|V|\times |V|$ with $a_{i,j}=1$ if and only if $ij\in E$, and $a_{i,j}=0$ otherwise. 
Since $A_G$ is real and symmetric, 
all eigenvalues of $A_G$ are real and can be reordered as $\lambda_1\ge \lambda_2\ge \cdots \ge \lambda_n$.  The eigenvalues of a graph $G$ are defined as the eigenvalues of the matrix $A_G$. 
The spectral radius $\lambda (G)$ is defined to be the maximum absolute value of eigenvalues of $G$. By the Perron--Frobenius theorem, 
the spectral radius of a nonnegative matrix is actually a largest eigenvalue. Furthermore, there exists a unit nonnegative eigenvector $\bm{x}\in \mathbb{R}^n$  corresponding to $\lambda (G)$. 
Such a vector is called the Perron--Frobenius eigenvector. For each vertex $v\in V(G)$, 
we write $x_v$ for the coordinate of $\bm{x}$ corresponding to $v$. 
In the setting of graphs, 
the eigen-equation $A_G\bm{x}= \lambda (G) \bm{x}$ can be interpreted as $\lambda(G) x_v = \sum_{u\in N(v)} x_u$, and the Rayleigh quotient implies 
$\lambda (G) = 2 \sum_{uv\in E} x_ux_v$. 
Here and in the rest of the paper, we denote by $\sum_{uv \in E}$ the sum over each edge in $E$ once. When $G$ is a bipartite graph with given partite sets $A$ and $B$, we additionally assume that $u \in A$ and $v \in B$.


\section{Preliminary results}\label{sec:prelim}
In this section, we collect a number of results that we will use throughout our proofs. The first one is a result of  Erd\H{o}s, Frankl and R\"{o}dl \cite{EFR1986} which allows us to delete $o(n^2)$ edges and transform an $F$-free graph to a $K_{r+1}$-free graph. 

\begin{lemma}[Erd\H{o}s--Frankl--R\"{o}dl \cite{EFR1986}] \label{lemEFR}
Let $F$ be a graph with $\chi (F)=r+1$. 
For any $\varepsilon >0$, 
there is $n_0$ such that 
if $n\ge n_0$ and $G$ is an $n$-vertex $F$-free graph, 
then we can remove at most $\varepsilon n^2$ edges from $G$ 
so that the remaining graph is $K_{r+1}$-free.  
\end{lemma}

The following lemma states that if $G$ is a minimal graph with a certain bound on the spectral radius, then the product of the Perron--Frobenius weights of each edge is at least $\Omega (m^{-1/2})$.  

\begin{lemma} \label{lem-product-entry}
Let $r\ge 2$, $\varepsilon>0$, $C>0$ and $\bm{x}=(x_v)_{v\in V(G)}$ be the unit Perron--Frobenius vector of $G$. 
Let $G$ be an $m$-edge graph with 
$  \lambda (G)> \sqrt{(1-\frac{1}{r} + \varepsilon )2m} + C$.  
 If for any $m'$-edge subgraph $G'\subseteq G$, $\lambda(G')\le 
\sqrt{(1-\frac{1}{r} + \varepsilon )2m'} + C$, then for any edge $uv \in E(G)$, we have  $x_ux_v > \frac{1}{8\sqrt{m}}.$
\end{lemma}

\begin{proof}
Suppose on the contrary that $x_ux_v \le 8^{-1}m^{-1/2}$ 
for some edge $uv\in E(G)$.  Then removing the edge 
$uv$ from $E(G)$, we get a subgraph $G'$ with $m'=m-1$ edges and 
\begin{align*}
\lambda (G') \ge  \left(2\sum_{ij \in E(G)} x_ix_j \right) - 2x_ux_v  \ge \lambda (G) - \frac{1}{4\sqrt{m}}  > 
 \sqrt{\Bigl( 1- \frac{1}{r} + 
 \varepsilon \Bigr)2m'} + C. 
 \end{align*}
 This contradicts with the assumption of $G$.  
\end{proof}


The following lemma says that if  $G$ is an $m$-edge graph with spectral radius bounded away from $\sqrt{m}$, then every coordinate of the Perron--Frobenius eigenvector of $G$ is small. 

\begin{lemma} \label{lem:max} 
Let $G$ be a graph with $m$ edges and $\bm{x}=(x_v)_{v\in V(G)}$ be the unit Perron--Frobenius eigenvector of $G$.  
 If $\lambda^2 (G) \ge (1+\delta)m$ where $0<\delta \le 0.79$, then 
\[ \max\{x_v : v\in V(G)\} < {\delta^{-4}m^{-1/4} }. \] 
\end{lemma}
\begin{proof}
Let $\lambda = \lambda (G)$ and 
$i\in V(G)$ be a vertex such that 
$x_i= \max\{x_v : v\in V(G)\}$. 
Suppose on the contrary that 
$x_i \ge \delta^{-4}m^{-1/4}$. Note that $\lambda \bm{x}= A_G \bm{x}$ gives 
\[ \lambda x_i = \sum_{j\in N(i)} x_j. \] 
Then we have 
\[  \lambda^2 x_i = \sum_{j\in N(i)} \lambda x_j= \sum_{j\in N(i)} 
\sum_{k\in N(j)} x_k = \sum_{k\in V(G)} d_{N(i)}(k) \cdot x_k , \]
where $d_{N(i)}(k)$ is the number of neighbors of $k$ in $N(i)$. Therefore, we get 
\begin{equation}\label{eq-eigen-sq}
\lambda^2 x_i \le \sum_{k\in V(G)} d_G(k)\cdot x_k = \sum_{jk \in E(G)}(x_j +x_k). 
\end{equation} 
For each edge $jk\in E(G)$, we call $jk$ \emph{good} if $\min\{ x_j ,x_k \} > \delta^{-2} m^{-1/4}$. Since 
$\sum_{v\in V} x_v^2 =1$, we get $\#\{ j\in V(G): x_j > \delta^{-2}m^{-1/4}\}  <  \delta^4 m^{1/2}$. Then 
the number of good edges of $G$ is at most $\delta^8m/2$. 
On the other hand, for each non-good edge 
$jk \in E(G)$, we have 
\[  x_j +x_k 
\le x_i + \delta^{-2} m^{-1/4} \le (1+\delta^2)x_i. \]  
Therefore, it follows from (\ref{eq-eigen-sq}) that 
\[  (1+\delta )m x_i \le \lambda^2 x_i \le 
2x_i \cdot \frac{1}{2} \delta^8m + (1+\delta^2) x_i \cdot m=(1+\delta^2 +\delta^8)mx_i . \]
Since $0<\delta <0.79$, we have 
$\delta > \delta^2 + \delta^8$, 
which leads to a contradiction. 
\end{proof}




\begin{lemma}
    \label{lem:remove-edges}
    If $G$ is a graph with $m$ edges, and $G'$ is the graph obtained from $G$ by removing at most $\gamma m$ edges, then $\lambda (G')> \lambda (G) - \sqrt{2\gamma m}$ and 
    $\lambda^2(G') > \lambda^2(G) - 4 \sqrt{\gamma} m.$
\end{lemma}

\begin{proof}
    By Rayleigh's formula, we have $\lambda(G') \geq \lambda(G) - \lambda(G \backslash G')$. Thus
    $$\lambda^2(G') \geq \lambda^2(G) - 2 \lambda(G \backslash G') \lambda(G).$$
    Using the bound $\lambda(G) < \sqrt{2m}$ and $\lambda(G \backslash G') < \sqrt{2 \gamma m}$, we get the desired result.
\end{proof}

The following well--known fact (see, e.g., \cite{BFP2008}) will be used in our argument. 

\begin{lemma} \label{lem-square-equal}
  Let $G$ be a bipartite graph with the partition $V(G)=A\sqcup B$, and $\binom{\bm{x}}{ \bm{y}}$ be an eigenvector of $G$ corresponding to $\lambda (G)$ with $\bm{x}$ indexed by $A$ and $\bm{y}$ indexed by $B$. Then we have $\lVert \bm{x}\rVert = \lVert \bm{y}\rVert$, i.e., 
  \[ \sum_{i\in A} x_i^2 = \sum_{j\in B} y_j^2. \]  
\end{lemma}

We provide a proof for the sake of completeness. 

\begin{proof}
  Note that $G$ is a bipartite graph. Then the adjacency matrix of $G$ is of the form:   
    \[  A_G=\begin{bmatrix} O & M \\ M^{T} & O \end{bmatrix}, \]
    where $M$ is a $(0,1)$-matrix of order $|A|\times |B|$ with 
    $m_{i,j}=1$ if the vertex $i\in A$ is adjacent to the vertex $j\in B$. As $\binom{\bm{x}}{\bm{y}}$ is the Perron--Frobenius eigenvector of $A_G$, we have 
    $M\bm{y} = \lambda(G) \bm{x}$ and 
    $M^{T}\bm{x}=\lambda(G) \bm{y}$. 
    Now, let $M=UDV$ be a singular value decomposition (see, e.g., \cite[p. 15]{Zhan2013}) of $M$, where $U,V$ are real orthogonal matrices, and $D$ is an $|A|\times |B|$ matrix with $(i,i)$-entries being the  singular values of $M$. Since $M\bm{y} = \lambda (G) \bm{x}$ and $U^TU=I$, we get  $DV\bm{y}=\lambda (G) U^{T} \bm{x}$, which implies that for each index $i$, 
    \[ D_{ii} \cdot (V\bm{y})_i =\lambda (G) \cdot (U^T\bm{x})_i.  \]
   Similarly, we obtain from $M^{T}\bm{x}= \lambda (G) \bm{y}$ that $D^TU^T \bm{x}= \lambda (G) V\bm{y}$. Then for each $i$, 
   \[ D_{ii} \cdot (U^T\bm{x})_i=\lambda (G)\cdot (V\bm{y})_i. \]
   Combining with the above two equalities, 
   we conclude that $|(V\bm{y})_i| = |(U^T\bm{x})_i|$, which yields  
    $\lVert V\bm{y}\rVert = \lVert U^T\bm{x}\rVert$. Since both $V$ and $U^T$ are orthogonal matrices, we get $\lVert \bm{y}\rVert = \lVert \bm{x}\rVert$. 
\end{proof}

We will also utilize the well-known Motzkin--Straus inequality \cite{MS1965}. 

\begin{lemma}[Motzkin--Straus \cite{MS1965}] \label{lem-MS} 
Let $r\ge 2$ be an integer and $G$ be a $K_{r+1}$-free graph on $n$ vertices. 
    If $\bm{x}=(x_i)_{i=1}^n$ is a non-negative real vector with 
$\sum_{i=1}^n x_i =1$, then 
\begin{equation*}
\sum_{ij \in E(G)} x_ix_j \le \frac{1}{2}\left( 1- \frac{1}{r} \right).
\end{equation*} 
The equality holds if and only if the subgraph induced by the vertices corresponding to non-zero entries of $\bm{x}$ is a complete $r$-partite graph such that the sum of the $x_i$'s in each part is the same. 
\end{lemma}

The following lemma will be used in the proof of \Cref{thm-common}. 

\begin{lemma}[Li--Liu--Zhang \cite{LLZ2024-book-qua}] \label{lem-Delta} 
 Let $G$ be a graph with $m$ edges and maximum degree $\Delta (G)$. 
If $\lambda (G)> \sqrt{m}$, then for sufficiently large $m$, the maximum degree $\Delta(G)$ of $G$ satisfies 
\[ \Delta (G)\le m/2 + m^{0.99}. \]
\end{lemma}









\section{Extensions on Erd\H{o}s--Stone--Simonovits' theorem}\label{sec:ESS}

\subsection{Proof of \Cref{conj-LLF}}  

Before proving \Cref{conj-LLF}, 
we present a short proof of \Cref{eq-spec-ESS} which is quite different from Nikiforov's proof in \cite{Niki2009cpc}. 
It illustrates the basic ideas in our proof of \Cref{conj-LLF}. 


\begin{proof}[{\bf Short proof of \Cref{thm-spec-ESS}}]
The lower bound can be witnessed by the Tur\'{a}n graph $T_{n,r}$ since $T_{n,r}$ is $F$-free and $\lambda (T_{n,r}) \ge (1-\frac{1}{r})n- \frac{r}{4n}$. 
Assume that $G$ is an $n$-vertex $F$-free graph. 
By Lemma \ref{lemEFR}, we can remove $o(n^2)$ edges from 
$G$ and get a $K_{r+1}$-free subgraph $G'$. Note that 
\[ \lambda (G\setminus G') <  
 \sqrt{2e(G\setminus G')} = o(n). \]   
 Since $G'$ is $K_{r+1}$-free, by Wilf's theorem, we have $\lambda (G') \le \left(1- \frac{1}{r} \right) n$. 
 Therefore, we obtain 
\[ \lambda (G) \le \lambda (G') + \lambda (G\setminus G') 
\le \left( 1- \frac{1}{r} \right)n + o(n), \]
which completes the proof.  
\end{proof}



Now we prove \Cref{conj-LLF}. 
In our proof of \Cref{thm-spec-ESS}, we remove $o(n^2)$ edges from $G$ 
to make it $K_{r+1}$-free, and then apply Wilf's bound. 
For \Cref{conj-LLF}, we need to show that the removal of these  edges changes the spectral radius by $o(m^{1/2})$, which is true in the dense case when $m=\Omega (n^2)$. Our key idea is to make a reduction to the case of dense graphs using Lemma \ref{lem:max}, which 
gives an upper bound on $\max_{v\in V(G)} \{x_v\}$, where $\bm{x}\in \mathbb{R}^n$ is the unit Perron--Frobenius eigenvector of $G$. 
We note that the assumptions of Lemma \ref{lem:max} do not directly imply that $G$ is dense. For example, $G$ can be the disjoint union of a clique $K_n$ with arbitrarily many isolated vertices. We will address this obstacle by a vertex deletion argument. 


\begin{proof}[{\bf Proof of \Cref{conj-LLF}}]
By Lemma \ref{lemEFR}, for each $\varepsilon > 0$, there exists some $N(\varepsilon) > 0$ such that any $F$-free graph on $n \geq N(\varepsilon)$ vertices can be made $K_{r + 1}$-free by removing at most $\varepsilon^{20} n^2$ edges. Assume that $G$ is an $F$-free graph with $m$ edges. 
It suffices to prove that for any $\varepsilon  \in (0,1/3)$, we have
$$\lambda(G)\le 
\sqrt{\Bigl(1-\frac{1}{r} + \varepsilon \Bigr)2m} + N(\varepsilon).$$ 
Suppose on the contrary that 
$G$ is an $F$-free graph with $m$ edges and 
\begin{equation} \label{eq-assume}
    \lambda (G)> \sqrt{\left(1-\frac{1}{r} + \varepsilon\right)2m} + N(\varepsilon).
\end{equation}
Among such counterexamples, we choose a graph $G$ such that $|V(G)| + |E(G)|$ is minimal. Denote $|V(G)|=n$. It is clear that $G$ has no isolated vertices. As removing any edge does not create a counterexample due to the minimality of $G$, 
then for any $m'$-edge subgraph $G' \subseteq G$, we have 
$\lambda (G')\le \sqrt{(1- \frac{1}{r} + \varepsilon)2m'} +N(\varepsilon)$. 
Using Lemma \ref{lem-product-entry},  
 for any edge $uv \in E(G)$, we have 
$$ x_ux_v > \frac{1}{8\sqrt{m}}.$$ 
Furthermore, by Lemma \ref{lem:max}, 
we have 
\[ \max \{x_v: v\in V(G)\}  \leq (2\varepsilon)^{-4}m^{-1/4}. \] 
These two bounds above, together with the fact that $G$ has no isolated vertex, imply that
$$\min \{ x_v: v\in V(G)\} > 2\varepsilon^4 m^{-1/4}.$$
As $\bm{x} \in \mathbb{R}^n$ is a unit vector, 
it follows that 
\begin{equation}
    \label{eq-dense}
    n \leq 4^{-1}\varepsilon^{-8} m^{1/2}.
\end{equation}
On the other hand, we get from (\ref{eq-assume}) that $n \geq \lambda(G) \geq N(\varepsilon)$. 
By \Cref{lemEFR}, we can remove at most $\varepsilon^{20} n^2 $ edges from $G$ to get a $K_{r+1}$-free subgraph $G'$ with $m'$ edges. Note that $\varepsilon^{20} n^2 < \varepsilon^4 m$ by (\ref{eq-dense}). 
Now we use the Nikiforov bound (\ref{eq-Niki-2002-cpc}) and we get 
\[  \lambda (G') \le 
\sqrt{\Bigl(1-\frac{1}{r}\Bigr)2m'}. \]
Thus, we conclude that  
\begin{equation*}\label{eq:removal}
    \lambda(G) \le \lambda(G') + \lambda(G\setminus G') < \sqrt{\Bigl(1-\frac{1}{r}\Bigr)2m'} + \sqrt{2 \varepsilon^4 m}, 
\end{equation*}
which contradicts with (\ref{eq-assume}). This completes the proof. 
\end{proof}



\subsection{Proof of \Cref{thm:walk}}  

To prove~\Cref{thm:walk}, we first show a generalization of \Cref{lem:max} for walks. 
Recall that an $\ell$-walk is a walk with $\ell$ vertices. 
For two vertices $i, j$ in a graph $G$, let $w_\ell(i)$ denote the number of $\ell$-walks starting from $i$, and let $w_\ell(i, j)$ denote the number of $\ell$-walks from $i$ to $j$.

\begin{lemma} \label{lem:max-II}
Let $G$ be a graph and $\bm{x}=(x_j)_{j\in V(G)}$ be its unit Perron--Frobenius eigenvector. If $\lambda^{\ell} (G) \ge \left(\frac{1}{2} + \delta\right) \cdot w_{\ell}(G)$ for some $\delta \in (0, 0.1)$, then 
\[ \max\{x_j : j \in V(G)\} < \delta^{-4}\lambda(G)^{-1/2}. \] 
\end{lemma}
\begin{proof}
We denote  $\lambda = \lambda(G)$ and let $i \in V(G)$ be a vertex with $x_i = \max\{x_j : j\in V(G)\}$. 
Assume for the sake of contradiction that 
$x_i \ge \delta^{-4}\lambda^{-1/2}$. We use the identity
$$\lambda^{\ell} x_i = \sum_{k \in V(G)} A_G^{\ell}(i,k) \cdot x_k = 
 \sum_{k \in V(G)} w_{\ell + 1}(i, k) \cdot x_k.$$ 
Note that each $(\ell +1)$-walk starting from the vertex $i$ 
is determined by the $\ell$-walk obtained by removing $i$. So we have
$$\lambda^{\ell} x_i \leq \sum_{j, k \in V(G)} w_\ell(j, k) \cdot x_k = \frac{1}{2} 
\sum_{j, k \in V(G)} w_\ell(j, k)\cdot  (x_j + x_k).$$
Let $B$ be a vertex subset of $G$ defined as below: 
\[ B := \left\{j\in V(G):  x_j > 
\delta^{-2} \lambda^{-1/2} \right\}. \] 
As $\bm{x}$ is unit, i.e., $\sum_{j\in V(G)} x_j^2 = 1$, we have 
\[ \abs{B} < \delta^4 \lambda.\] 
We call an $\ell$-walk $w$ from $j$ to $k$ \emph{good} if both $j$ and $k$ lie in $B$, and \emph{bad} otherwise.
Applying the Rayleigh's formula, the number of good $\ell$-walks in $G$ is given by
$$\sum_{j, k \in B} w_\ell(j, k) = \mathbf{1}_{B}^T A_G^{\ell - 1}  \mathbf{1}_{B} \leq \lambda^{\ell - 1} \mathbf{1}_{B}^T \mathbf{1}_{B} = \lambda^{\ell - 1} |B| \leq \delta^4 \lambda^{\ell}.$$
Thus we have
\begin{equation} \label{upper-bad}
    \frac{1}{2} \sum_{j, k \in B} w_\ell(j, k) \cdot (x_j + x_k) \leq \delta^4 \lambda^{\ell} x_i.
\end{equation}
On the other hand, for any bad $\ell$-walk from $j$ to $k$, we have 
\[ x_j+x_k\le x_i+\delta^{-2}\lambda^{-1/2}\le(1+\delta^2)x_i. \]
Thus, the contribution from bad walks is bounded by 
\begin{equation} \label{upper-not-bad}
    \frac{1}{2}
\sum_{\substack{j, k \in V(G)\\ \{j, k\} \not\subseteq B}} w_\ell(j, k) \cdot (x_j + x_k) \leq \frac{1}{2} w_{\ell}(G) \cdot (1 + \delta^2) x_i \leq \frac{1 + \delta^2}{1 + \delta} \lambda^{\ell} x_i,
\end{equation}
where the last inequality follows by the assumption  $\lambda^{\ell}\ge \left(\frac{1}{2} + \delta\right) \cdot w_{\ell}(G)$. 

Summing the two inequalities (\ref{upper-bad}) and (\ref{upper-not-bad}), we conclude that
$$\frac{1}{2}\sum_{j, k \in V(G)} w_\ell(j, k) \cdot (x_j + x_k) \leq \left(\frac{1 + \delta^2}{1 + \delta} + \delta^4\right) \lambda^{\ell} x_i < \lambda^{\ell} x_i,$$
which leads to a contradiction.
\end{proof}

An additional lemma we need is the following. 

\begin{lemma}
    \label{lem:walks-lower}
    Let $G$ be a graph and $\bm{x}=(x_j)_{j\in V(G)}$ be its unit Perron--Frobenius eigenvector. Then for any vertex $i$ in $G$, we have
    $$w_\ell(i) \geq \frac{\lambda^{\ell - 1}(G) x_i}{\max_{j \in V(G)} \{x_j\}}.$$
\end{lemma}
\begin{proof}
    Let $\lambda=\lambda(G)$ and $\mathbf{1}_i$ be the vector with $1$ at the $i$-th entry and $0$ everywhere else. Then
    $$\mathbf{1}_i^T A_G^{\ell - 1} \bm{x} = \lambda^{\ell - 1}  \mathbf{1}_i^T \bm{x}  = \lambda^{\ell - 1} x_i.$$
    On the other hand, it follows that 
    $$ \mathbf{1}_i^T A_G^{\ell - 1} \bm{x}   
   = \sum_{j \in V(G)}w_\ell(i, j) \cdot x_j \leq 
     \sum_{j \in V(G)}w_\ell(i, j) \cdot \max_{j \in V(G)} \{ x_j \} = w_\ell(i) \cdot \max_{j \in V(G)} \{ x_j \},$$
    as desired.
\end{proof}

We can now prove \Cref{thm:walk}.

\begin{proof}[{\bf Proof of \Cref{thm:walk}}]
By Lemma \ref{lemEFR}, for each $\varepsilon > 0$, there exists some $N(\varepsilon) > 0$ such that any $F$-free graph on $n \geq N(\varepsilon)$ vertices can be made $K_{r+ 1}$-free by removing at most $\varepsilon^{20}\ell^{-8} n^2/256$ edges. Assume that $G$ is an $F$-free graph with $m$ edges. 
It suffices to prove that, for any $0< \varepsilon <0.01$, we have
$$\lambda^{\ell}  (G)\le 
\left(1-\frac{1}{r} + \varepsilon\right) \cdot w_{\ell}(G) + N(\varepsilon)^{\ell}.$$
Suppose  on the contrary that there exists a graph $G$ with
\begin{equation} \label{eq-contra}
\lambda^{\ell} (G) > 
\left(1-\frac{1}{r} + \varepsilon\right) \cdot w_{\ell}(G) + N(\varepsilon)^{\ell}.
\end{equation}
Among such $G$, we choose a graph $G$ such that $|V(G)| + |E(G)|$ is minimal. Denote $V = V(G)$. It is clear that $G$ does not have an isolated vertex. Let $e = \{i, j\}$ be any edge in $G$, and let $G' = G \backslash \{e\}$ be the subgraph of $G$ obtained by removing $e$. 
By the minimality of $G$, we have
$$\lambda^{\ell}(G') \leq \left(1-\frac{1}{r} + \varepsilon\right) \cdot w_{\ell}(G') + N(\varepsilon)^{\ell}.$$
Therefore, we observe that 
$$\lambda^{\ell}(G) - \lambda^{\ell}(G') > \left(1- \frac{1}{r} + \varepsilon\right) \cdot (w_{\ell}(G) - w_{\ell}(G')).$$ 
On the one hand, we recall that 
$ \lambda(G')\ge \lambda(G) - 2 x_i x_j$. Recall the elementary inequality that $a^{\ell}-b^{\ell} 
\le \ell a^{\ell -1}(a-b)$ holds for any numbers $a\ge b>0$. Then 
\[  \lambda^{\ell}(G)- \lambda^{\ell}(G') \le \ell \cdot \lambda^{\ell -1}(G) \cdot 2x_ix_j.  \]
On the other hand, since any $\ell$-walk in $G$ starting with the edge $ji$ in $G$ no longer appears in $G'$, and there are $w_{\ell - 1}(i)$ such walks, we also have
$$w_{\ell}(G) - w_{\ell}(G') \geq w_{\ell - 1}(i).$$  
Therefore, we have
$$\ell  \cdot \lambda^{\ell - 1}(G)\cdot 2x_i x_j \geq 
\left(1- \frac{1}{r} + \varepsilon\right) \cdot w_{\ell - 1}(i).$$
Invoking \Cref{lem:walks-lower}, we have
$$w_{\ell - 1}(i) \geq \frac{\lambda^{\ell - 2}(G) x_i}{\max_{k \in V} \{x_k\}}.$$
So we conclude that
$$x_j \cdot \max_{k \in V} \{x_k\} \geq \frac{1}{4\ell \lambda(G)}.$$
As each vertex $j$ in $G$ is contained in at least one edge, this inequality holds for any vertex $j$.

By \Cref{lem:max-II}, we have
$\max \{x_k:k \in V\} < \varepsilon^{-4} \lambda(G)^{-1/2}$. 
Then for any $j \in V$, 
$$x_j \geq \frac{\varepsilon^4}{4\ell} \lambda(G)^{-1/2},$$
which implies
$$\abs{V} \leq 16\ell^2 \varepsilon^{-8} \lambda(G).$$
As $|V|\ge\lambda(G)\ge N(\varepsilon)$, by the removal Lemma, we can remove at most $\varepsilon^{20}\ell^{-8} \abs{V}^2/256 \leq \varepsilon^4 \ell^{-4}\lambda^2(G)$ edges from $G$ to obtain a $K_{r+1}$-free subgraph $G'$. 
Applying \Cref{thm-Niki-2006} on $G'$ gives 
$$\lambda(G') \leq \left(\Bigl(1-\frac{1}{r}\Bigr) w_{\ell}(G')\right)^{1/\ell} \leq \left(\Bigl(1-\frac{1}{r}\Bigr) w_{\ell}(G)\right)^{1/\ell}.$$
So we conclude that
\[ \lambda(G) \le \lambda(G') + \lambda(G\setminus G') <\left(\Bigl(1-\frac{1}{r}\Bigr) w_{\ell}(G)\right)^{1/\ell} + \sqrt{2\varepsilon^4 \ell^{-4} \lambda^2(G)}. \]  
Simplifying, we obtain
\[  \lambda^{\ell} (G) < \left(1- \sqrt{2} \varepsilon^2\ell^{-2} \right)^{-\ell} 
 \left(1-\frac{1}{r}\right) w_{\ell}(G), \]
which together with (\ref{eq-contra}) gives 
\[  \left(1- \frac{1}{r} \right)  > \left(1- \sqrt{2} \varepsilon^2\ell^{-2} \right)^{\ell} \left(1-\frac{1}{r} + \varepsilon\right) \ge 
\left(1- \sqrt{2}\varepsilon^2\ell^{-1}\right) \left(1-\frac{1}{r} + \varepsilon\right), \]
where the second inequality follows by Bernoulli's inequality. 
So we get $ \sqrt{2}\varepsilon \ell^{-1}(1-\frac{1}{r}+\varepsilon)>1$, which leads to a contradiction since $\varepsilon< 0.01$.  
The proof is complete. 
\end{proof}









\section{The edge-spectral stability theorem}
\label{sec:edge-stability} 


Our proof of~\Cref{thm:nikiforov-stability} for general graphs $F$ reduces to the case when $F$ is a clique. We first prove the clique case when $F$ is a triangle in~\Cref{sec:F-triangle},  and then the case when $F=K_{r+1}$ with $r\ge 3$  in~\Cref{sec:F-clique}. Finally, we settle the general case in~\Cref{sec:F-general}.


\subsection{The case of triangles}\label{sec:F-triangle}

As a starting point for the whole argument, let us prove \Cref{thm:nikiforov-stability} when $F$ is the triangle $K_3$. 
For the convenience of readers, we extract this case as the following theorem. 

\begin{theorem} \label{thm-stability-triangle}
    For every $\varepsilon\in (0,0.01)$, 
    there exists $\delta = \delta (\varepsilon) >0$ such that if $G$ is a triangle-free graph with $m$ edges and 
    \[ \lambda (G)\ge (1- \delta )\sqrt{m}, \] 
    then there exist disjoint vertex subsets $A,B \subseteq V(G)$ such that $d(G,K_{A,B})\le \varepsilon m$.
\end{theorem}

Recall that the stability of the Mantel theorem says that 
if $G$ is an $n$-vertex triangle-free graph with 
$e(G)\ge n^2/4 -q$, then 
$G$ can be made bipartite by removing at most $q$ edges. Next, we show a spectral version of this statement.  

\begin{lemma} \label{lem-bipartite}
For every $\varepsilon >0$, 
if $G$ is an $m$-edge triangle-free graph with
\[ \lambda (G) \ge 
(1- \varepsilon /2 ) \sqrt{m}, \] 
then $G$ can be made bipartite by removing at most 
$\varepsilon m$ edges. 
\end{lemma}

\begin{proof}  
Let $u\in V(G)$ be a vertex such that 
$x_u= \max\{x_v: v\in V(G)\}$. 
We denote $U:=N(u)$ and $W:=V(G)\setminus (N(u)\cup \{u\})$.  
Then 
\[  \lambda^2 x_u =\sum_{j\in N(u)} \sum_{k\in N(j)} x_k 
= d(u)x_u + \sum_{v\in U} d_U(v)x_v + 
\sum_{w\in W}d_U(w) x_w. \]
Since $G$ is triangle-free, we have $d_U(v) = 0$ whenever $v \in U$. Thus we have
$$(1- \varepsilon )mx_u < \lambda^2 x_u= d(u)x_u + \sum_{w\in W}d_U(w) x_u = (d(u) + e(U, W)) x_u.$$
It follows that 
$$e(W) = 
m- d(u) - e(U,W) <  \varepsilon m.$$ 
After removing all edges within 
$W$, since there are at most $\varepsilon m$ edges in $W$, we obtain the desired bipartite graph with two vertex parts $\{u\}\cup W$ and $U$.  
\end{proof}

We note that this gives the following concise vertex-spectral stability theorem.


\begin{cor}  \label{cor-bipartite-vertex-spectral}
For any $0\le \varepsilon \le 1$ and $n\ge 4/\varepsilon$, if $G$ is an $n$-vertex triangle-free graph with 
\[   \lambda (G) \ge 
(1- \varepsilon/2 )\cdot  \lambda (T_{n,2}), \]  then $G$ can be made bipartite by removing at most $\varepsilon n^2$ edges. 
\end{cor}

\begin{proof}
    The Rayleigh formula implies $\lambda (G)\ge \frac{2m}{n}$, which together with the assumption gives  
    \[ \lambda^2(G)\ge \frac{2m}{n} \lambda (G) \ge 
    \frac{2m}{n} \left(1- \frac{\varepsilon}{2}\right) 
    \lambda (T_{n,2}).\]  
    Note that $\lambda (T_{n,2}) \ge \delta (T_{n,2}) =\lfloor \frac{n}{2}\rfloor \ge \frac{n-1}{2}$. 
    To apply Lemma \ref{lem-bipartite}, it suffices to show that 
    \[ \frac{2m}{n}\left(1- \frac{\varepsilon}{2}\right) \frac{n-1}{2} \ge 
    \left(1-\frac{\varepsilon}{2}\right)^2 m, \] 
which is equivalent to showing that 
\[ \left(1- \frac{\varepsilon}{2}\right) {(n-1)} 
\ge 
\left(1- \varepsilon+ \frac{\varepsilon^2}{4} \right)n. \] 
    The inequality holds by $0\le \varepsilon \le 1$ and 
    $n\ge 4/\varepsilon$, and then  
    $(\frac{\varepsilon}{2} -\frac{\varepsilon^2}{4}) n \ge \frac{1}{4}\varepsilon n > 1- \frac{\varepsilon}{2}$. 
\end{proof}

In what follows, we give another proof of Corollary \ref{cor-bipartite-vertex-spectral}. 

\begin{proof}[Second proof]
Since $G$ is a triangle-free graph with $m$ edges, we know that $\lambda (G)\le \sqrt{m}$. 
    First of all, we claim that $e(G)\ge \lfloor n^2/4 \rfloor  -\varepsilon n^2$. 
    Otherwise, if $e(G)< \lfloor n^2/4 \rfloor  -\varepsilon n^2$, then 
    \[ \lambda (G)\le \sqrt{m} 
    < \sqrt{\big\lfloor {n^2}/{4} \big\rfloor  -\varepsilon n^2} \le \left(1- \frac{\varepsilon}{2}\right)\lambda (T_{n,2}), \]  
    where the last inequality holds since $\lambda (T_{n,2}) = \sqrt{\lfloor n^2/4\rfloor}$ and  $\varepsilon \ge 0$. This contradicts with the assumption. Thus, we conclude that $e(G)\ge \lfloor n^2/4\rfloor - \varepsilon n^2$. 

    Let $A$ be an independent set of maximum size in $G$. 
Since $G$ is triangle-free, $N(v)$ is an independent set for every $v\in V(G)$, which implies $\Delta(G) \le |A|$. 
Setting $B:=V(G)\setminus A$, we get 
\[ \sum_{v\in B} d(v) \le |B| \Delta (G) \le |B| |A| 
\le \left\lfloor  {n^2}/{4} \right\rfloor.  \]  
On the other hand, it follows that 
\[  \sum_{v\in B} d(v) =e(A,B) + 2e(B) = e(G) + e(B).  \]
Since $e(G)\ge \lfloor {n^2}/{4}\rfloor - \varepsilon n^2$, we get $e(B)\le \varepsilon n^2$. 
Removing all edges within $B$, we can obtain a  bipartite subgraph of $G$ on the vertex parts $A$ and $B$, as needed.   
\end{proof}


To complete the proof of~\Cref{thm-stability-triangle}, 
we show in the rest of this subsection that the resulting bipartite subgraph obtained in Lemma \ref{lem-bipartite} 
is close to a complete bipartite graph. 
Throughout, we write $\bm{x}$ for the unit Perron--Frobenius eigenvector of $G$.


\begin{lemma}
\label{lem:close-to-Kab}
For every $\varepsilon \in (0,0.01)$, 
let $\delta = \varepsilon^4/100$. 
If $G$ is a bipartite graph with 
$m$ edges and $\lambda (G) \ge 
(1- \delta ) \sqrt{m}$, then there exist disjoint vertex subsets $U, V\subseteq V(G)$ such that 
\[ d(G, K_{U, V}) \leq \varepsilon m. \]  
\end{lemma}

\begin{proof}
 Let $V(G)=A\sqcup B$ be a bipartition of $G$. Without loss of generality, we may assume that $A = \{1,2,\ldots ,a\}$ and $B = \{1,2,\ldots ,b\}$, and the unit Perron--Frobenius eigenvector of $G$ takes values $x_1 \geq \dots \geq x_a$ on $A$ and $y_1 \geq \dots \geq y_b$ on $B$.   
 It follows from Lemma \ref{lem-square-equal} that
    $$\sum_{i \in A} x_i^2 = \sum_{j \in B} y_j^2 = \frac{1}{2}.$$
    Therefore, we have
    $$\sum_{i \in A} \sum_{j \in B} x_i^2y_j^2 = \frac{1}{4}.$$
    On the other hand, the assumption gives 
    $$\sum_{ij \in E(G)} 2x_i y_j = \lambda(G) \geq (1 - \delta) \sqrt{m}.$$
    Combining the two inequalities, we obtain
    \begin{equation}
    \label{eq:Motzkin-Straus-Exploit-I}
    \sum_{ij \notin E(G)} x_i^2 y_j^2 + \sum_{ij \in E(G)} \left(x_i y_j - \frac{1}{2\sqrt{m}}\right)^2 \leq \delta/2,
    \end{equation}
    where $\sum_{ij \notin E(G)}$ goes over each pair $(i, j) \in (A \times B )\backslash E(G)$ exactly once. 

    \medskip 
    Let $r, s$ be the smallest integers such that
   $\sum_{i = 1}^r x_i^2 \geq \varepsilon^2 
    ~\text{and}~  
    \sum_{j = 1}^s y_j^2 \geq \varepsilon^2$, respectively.   

\begin{claim}\label{cl:1}
We have $\frac{1}{3\sqrt{m}} < x_r y_s < \frac{1}{\sqrt{m}}$.
\end{claim}


    \begin{poc}
    It is easy to verify that $(x-\frac{1}{2\sqrt{m}})^2 \ge \frac{1}{4}x^2$ holds for every $x\in (-\infty, \frac{1}{3\sqrt{m}})\cup  (\frac{1}{\sqrt{m}}, +\infty)$. 
    If $x_r y_s \ge \frac{1}{\sqrt{m}}$, then 
    for any $i \le r$ and $j \le s$, we have $x_i y_j \geq \frac{1}{\sqrt{m}}$. This implies
    $$\min\left\{ x_i^2 y_j^2, \left(x_i y_j- \frac{1}{2\sqrt{m}} \right)^2 \right\} \geq \frac{1}{4} x_i^2 y_j^2.$$
    Summing over $i\le r$ and $j\le s$, we have
    $$\sum_{i=1}^r \sum_{j=1}^s \min\left\{ x_i^2 y_j^2, \left(x_i y_j - \frac{1}{2\sqrt{m}}\right)^2 \right\} \geq \frac{1}{4} \left(\sum_{i =1}^r x_i^2 \right) \left( \sum_{j=1}^s y_j^2 \right) \geq \frac{1}{4} \varepsilon^4$$
    which contradicts with  (\ref{eq:Motzkin-Straus-Exploit-I}). If $x_r y_s \leq \frac{1}{4\sqrt{m}}$, then $x_iy_j \le \frac{1}{4\sqrt{m}}$ 
    for $i \ge r$ and $j \ge s$. 
    We can finish by an analogous argument that replaces $i \le r$ and $j \le s$ with $i \geq r$ and $j \geq s$.
    \end{poc}

     Let $U$ and $V$ be sets of vertices of $G$ defined as  
    \[ 
    U := \left\{u\in A: x_u \geq \frac{x_r}{ 4} \right\} 
    \quad~\text{and}~\quad 
    V := \left\{v \in B: y_v \geq \frac{y_s}{4} \right\}. \]
    In the rest of the proof, we use (\ref{eq:Motzkin-Straus-Exploit-I}) to show that $G$ has small edit distance with $K_{U, V}$. 
    
    First of all, we show that there are few  vertex pairs in $U \times V$ that are non-edges in $G$. 

\begin{claim}\label{cl:2}
We have $|E(K_{U, V}) \backslash E(G)| \leq \frac{1}{3} \varepsilon m $.
 \end{claim}  
  
    \begin{poc}
    We need to show that $G$ misses at most $\frac{1}{3} \varepsilon m$ edges between $U$ 
    and $V$.  By~\Cref{cl:1}, for any $(i, j) \in (U\times V) \backslash E(G)$, we have 
    $$x_i y_j \geq \frac{1}{16} x_r y_s  \geq  \frac{1}{48 \sqrt{m}}.$$ 
    Using (\ref{eq:Motzkin-Straus-Exploit-I}), we obtain
        $$|E(K_{U, V}) \backslash E(G)| \leq 
        (48\sqrt{m})^2 \cdot \sum_{ij \notin E(G)} x_i^2 y_j^2
        \leq 
        (48 \sqrt{m})^{2} \cdot \frac{\delta}{2} 
        \leq 1152 \delta m \leq \frac{1}{3} \varepsilon m ,$$
        as desired.
    \end{poc}

    Moreover, we prove that there are few edges in $G[\bar{U},B]$ and 
    $G[A,\bar{V}]$. 

\begin{claim}\label{cl:3}
    We have $e(E(G) \cap (\bar{U} \times B)) \leq \frac{1}{3} \varepsilon m $, and symmetrically $e(E(G) \cap (A \times \bar{V})) \leq \frac{1}{3} \varepsilon m $.
\end{claim} 
    
    \begin{poc}
        Let $E_1$ be the set of edges $ij$ in $E(G) \cap (\Bar{U} \times B)$ with $j \geq s$, and let $E_2$ be the set of edges $ij$ in $E(G) \cap (\Bar{U} \times B)$ with $j < s$. It suffices to show that $\max\{\abs{E_1}, \abs{E_2}\} \leq 
        \frac{1}{6} \varepsilon m $.

        For any $ij \in E_1$, by Claim \ref{cl:1}, we have $x_i y_j \leq \frac{1}{4} x_r y_s  \leq \frac{1}{4\sqrt{m}}$. Then together with (\ref{eq:Motzkin-Straus-Exploit-I}), we have 
        $$\abs{E_1} \cdot \frac{1}{16 m} \leq \sum_{ij \in E_1} \left(x_i y_j - \frac{1}{2\sqrt{m}} \right)^2 \leq \frac{1}{2}\delta , $$
        which implies $\abs{E_1} \leq 8 \delta m < \frac{1}{6} \varepsilon m $, as needed. 

        To bound $|E_2|$, we use the Cauchy--Schwarz inequality to obtain
        \begin{align}
       \notag \sum_{ij \in E(G) \backslash E_2} 2x_i y_j  & \leq 2 \sqrt{\sum_{ij \in E(G) \backslash E_2} x_i^2 y_j^2} \cdot \sqrt{ |E(G) \backslash E_2| } \\ 
        & \leq 2 \sqrt{\sum_{i \in A, j \in B} x_i^2 y_j^2} \cdot \sqrt{|E(G) \backslash E_2|} \notag \\
        &= \sqrt{m - |E_2| }. \label{eq-notin-E2}
        \end{align}
        By the definition of $s$, we have $\sum_{j < s} y_j^2 < \varepsilon^2$. Observe that 
        $$\sum_{ij \in E_2} x_i^2 y_j^2 \leq \left(\sum_{i \in A} x_i^2\right) \left( \sum_{j < s} y_j^2 \right) < \frac{\varepsilon^2}{2}.$$
        Applying the Cauchy--Schwarz inequality, we obtain
        \begin{align} \label{eq-in-E2}
        \sum_{ij \in E_2} 2x_i y_j \leq 2 \sqrt{\sum_{ij \in E_2} x_i^2 y_j^2} \cdot \sqrt{\abs{E_2}} \leq \sqrt{2 \varepsilon^2 \abs{E_2}}.
        \end{align}
        Combining (\ref{eq-notin-E2}) and (\ref{eq-in-E2}), we conclude that 
        $$ 
        (1 - \delta) \sqrt{m} \le 
        \lambda(G) = \sum_{ij \in E(G)} 2x_i y_j 
        \le \sqrt{m - \abs{E_2}} + \sqrt{2 \varepsilon^2 \abs{E_2}}.$$
        By the Cauchy--Schwarz inequality again, we have
        $$\sqrt{m - \abs{E_2}} + \sqrt{2 \varepsilon^2 \abs{E_2}} \leq \sqrt{(1 + 4\varepsilon^2) \left(m - \abs{E_2} + \frac{1}{2}\abs{E_2}\right)}.$$
        Therefore, we obtain
        $$ (1 - \delta)^2 m\le 
        (1 + 4\varepsilon^2) \left(m - \abs{E_2} + \frac{1}{2}\abs{E_2} \right), $$
        which implies that $\abs{E_2} \leq \frac{1}{6} \varepsilon m $ by our choice of $\delta = \varepsilon^4 / 100$ and $\varepsilon \in (0, 0.01)$.
    \end{poc}
    Combining~\Cref{cl:2} and~\Cref{cl:3}, we conclude that 
    $$d(G, K_{U, V}) \leq e(K_{U, V} \backslash G) + e(E(G) \cap (\bar{U} \times B)) + e(E(G) \cap (A \times \bar{V})) \leq \varepsilon m, $$ 
    as desired.
\end{proof}

We can now prove \Cref{thm-stability-triangle}. 


\begin{proof}[{\bf Proof of \Cref{thm-stability-triangle}}]
We take $\delta = \frac{1}{4}\varepsilon^{36}$. 
Assume that $G$ is a triangle-free graph with $m$ edges and $\lambda (G)\ge (1- \frac{1}{4}\varepsilon^{36})\sqrt{m}$. 
Using Lemma \ref{lem-bipartite}, we obtain a bipartite subgraph $G'\subseteq G$ such that $G'$ has $m'\ge m - \frac{1}{2}\varepsilon^{36} m $ edges. 
Note that the Rayleigh's formula gives 
$$\lambda (G)\le \lambda (G') + \lambda (G\setminus G') \le \lambda (G') + \sqrt{2(m-m')}\le \lambda (G') + \varepsilon^6\sqrt{m}.$$ 
Then $\lambda (G')\ge (1-\delta)\sqrt{m} - \varepsilon^6 \sqrt{m} \ge (1-2\varepsilon^6) \sqrt{m'}> (1-\frac{\varepsilon^4}{1600})\sqrt{m'}$. Now, we can apply~\Cref{lem:close-to-Kab} to $G'$ to obtain two disjoint vertex sets $U,V \subseteq V(G')$ such that $d(G',K_{U,V})\le \frac{\varepsilon}{2}m$. Thus, we conclude that $d(G,K_{U,V})\le d(G,G') + d(G',K_{U,V})\le \frac{1}{2}\varepsilon^{36} m + \frac{1}{2}\varepsilon m< \varepsilon m$, as desired. 
\end{proof}


\subsection{The case of cliques}\label{sec:F-clique}

In this subsection, we prove \Cref{thm:nikiforov-stability} for larger forbidden cliques $F = K_{r + 1}$ with $r\ge 3$. 
We rephrase this case as the following theorem. 

\begin{theorem} 
\label{thm-stability-clique-4}
    For every $r\ge 3$ and $\varepsilon>0$, there exist $m_0$ and $\delta >0$ such that if $G$ is a $K_{r+1}$-free graph with $m\ge m_0$ edges and 
    \[ \lambda^2(G)\ge 
    \left( 1-\frac{1}{r}-\delta \right) 2m, \] 
    then there exists a vertex set $C\subseteq V(G)$ such that $d(G,T_{C,r})\le \varepsilon m$. 
\end{theorem}
 

We provide two different proofs of \Cref{thm-stability-clique-4}. 
The first proof is based on the Motzkin--Straus inequality \Cref{lem-MS}, and then reduce the edge-spectral stability in \Cref{thm-stability-clique-4} to the vertex-spectral stability in \Cref{thm-niki-stability}. 
The second alternative proof applies a recent result involving the stability of the Kruskal--Katona type result for hypergraphs. We present the first proof here and postpone the second proof to the Appendix \ref{sec-A}. 

First of all, we exploit the slackness in Nikiforov's application of the Motzkin--Straus inequality. Using this slackness, we show that $G$ contains an induced subgraph $G'$ such that $d(G,G')=o(m)$ 
and $\lambda (G')= (1-\frac{1}{r}- o(1))|G'|$. This puts us in a position to use \Cref{thm-niki-stability}.

\begin{lemma}
\label{lem:nikiforov-stability-clique}  
For every $\varepsilon > 0$ and $r \geq 3$, 
there exist $\delta>0$ and  $m_0$ such that the following holds. If $G$ is a $K_{r + 1}$-free graph with $m\ge m_0$ edges and $\lambda^2 (G)\ge (1- \frac{1}{r} - \delta) 2m$, then $G$ has an induced subgraph $G'$ such that $\abs{V(G')} \leq (1 + \varepsilon) \sqrt{2m / (1 - 1/r)}$ and $d(G, G') \leq \varepsilon m$.
\end{lemma}
\begin{proof}
    In the following argument, we can take $\delta = (0.1\varepsilon)^8$.  
    For simplicity, we write $V$ and $E$ for the set of vertices and the set of edges of $G$, respectively, and 
    let $\bm{x}=(x_i)_{i\in V}$ be the unit Perron--Frobenious eigenvector of $G$. 
    By Lemma \ref{lem-MS}, we have
    \begin{equation*}\label{eq:MS}
        \sum_{ij \in E} x_i^2x_j^2 \leq \frac{r - 1}{2r}.
    \end{equation*}
    On the other hand, the Rayleigh quotient gives 
    $$\sum_{ij \in E} 2x_i x_j = \lambda(G) \geq \sqrt{\Big(\frac{r-1}{r} - \delta \Big) 2m}.$$
    Setting $\alpha :=(\frac{r-1}{r}\cdot \frac{1}{2m} )^{1/4}$, a linear combination of the above inequalities yields that
    \begin{equation} \label{eq-expect}
    \sum_{ij \in E} \left(x_i x_j - \alpha^2\right)^2 \leq \frac{r - 1}{r} - \sqrt{\frac{r-1}{r} \cdot \Big(\frac{r - 1}{r}-\delta\Big)} < \delta.
    \end{equation}
    We say that an edge $ij \in E$ is \emph{bad} if $\abs{x_i x_j - \alpha^2} > \delta^{1/4} \alpha^2$, and \emph{good} otherwise. Let $G_b$ be the subgraph of bad edges, and $G_g$ be the subgraph of good edges. 
    For simplicity, we write $E_b$ and $E_g$ for the edge sets of $G_b$ and $G_g$, respectively. 
    From (\ref{eq-expect}), we have 
    \begin{equation} \label{upper-Eb}
         |E_b| \le \frac{\delta}{\delta^{1/2}\alpha^4} \leq 4\delta^{1/2} m.
    \end{equation}
    Using this fact, we now show that almost all vertices $i$ satisfy $x_i \approx \alpha$.

    For each positive integer $t \geq 1$, let 
    $$\Gamma_t:=\left\{i \in V: |x_i - \alpha| \geq 2\delta^{1/4} \alpha t \right\}.$$ 
    Observe that $\Gamma_t \subset \Gamma_{t - 1}$. We proceed via the following claims.

    The first one follows by the definition of $G_g$.
    
    \begin{claim}\label{cl:obs1}
    The graph $G_g[\Gamma_1]$ is a bipartite graph on the vertex parts 
   \[ \Gamma_1^+ := \{i \in V: x_i \geq (1 + 2\delta^{1/4}) \alpha\} \]
    and 
    \[ \Gamma_1^- := \{i \in V: x_i \leq (1 - 2\delta^{1/4}) \alpha\}. \]
    \end{claim}

    \begin{claim}\label{cl:obs2}
    For each $t \leq \delta^{-1/8} / 2$, all the edges between $\Gamma_t$ and $V \backslash \Gamma_{t - 1}$ are bad. 
    \end{claim}
    
    \begin{poc}
    Indeed, if $i \in \Gamma_t$ and $j \in V \backslash \Gamma_{t - 1}$, then
    \begin{align*}
    \abs{x_i x_j - \alpha^2} &= \abs{(x_i - \alpha)x_j + (x_j - \alpha)\alpha} \geq x_j\abs{x_i - \alpha} - \alpha 
    \abs{x_j - \alpha} \\
    &\geq (1 - 2(t - 1)\delta^{1/4}) \cdot 2t\delta^{1/4}\alpha^2  -
    2(t - 1)\delta^{1/4} \alpha^2 \\
    &> 2\delta^{1/4} \alpha^2 
    - 4t^2 \delta^{1/2} \alpha^2\geq \delta^{1/4} \alpha^2, 
    \end{align*}
    so the edge $ij$ is bad.
    \end{poc}

    We now consider the quantities 
    $$s_t := \sum_{i \in \Gamma_t} x_i^2.$$
    Our goal is to find an integer $\ell$ such that 
    $s_{\ell} \le O(\delta^{1/4})$, and then we will show that there are at most $\varepsilon m$ edges incident to the vertices of $\Gamma_{\ell}$ and that deleting all vertices of $\Gamma_{\ell}$ yields the desired subgraph $G'$.

    
    For every $t \le \delta^{-1/8} / 2$, as $\bm{x}$ is the Perron--Frobenius eigenvector, we have
    $$\lambda(G) s_t = \sum_{i \in \Gamma_t} \sum_{j \in N(i)} x_i x_j 
    =\sum_{i \in \Gamma_t} \sum_{\substack{j \in N(i)\\ j\in \Gamma_{t-1}}} x_ix_j + 
    \sum_{i \in \Gamma_t} \sum_{\substack{j \in N(i)\\ j\in V\setminus \Gamma_{t-1}}} x_ix_j.$$
    Note that $\Gamma_t \subset \Gamma_{t-1}$ for each $t\ge 1$. 
    By~\Cref{cl:obs2}, every edge summed over on the right hand side is either in $G_g[\Gamma_{t - 1}]$ or bad. Then 
    $$\lambda(G) s_t \leq \sum_{ij \in E_g[\Gamma_{t - 1}]} 2x_i x_j + \sum_{ij \in E_b} x_i x_j.$$
    By~\Cref{cl:obs1}, $G_g[\Gamma_{t - 1}]\subseteq G_g[\Gamma_{1}]$ is bipartite, so its spectral radius is at most $\sqrt{m}$. Thus, 
    $$\sum_{ij \in E_g[\Gamma_{t - 1}]} 2x_i x_j \leq \sqrt{m}\sum_{i \in \Gamma_{t - 1}} x_i^2 = \sqrt{m} s_{t - 1}.$$
    On the other hand, recall in (\ref{upper-Eb}) that  
    $|E_b|\le 4\delta^{1/2}m$. Then    
    $$\sum_{ij \in E_b} x_i x_j < \frac{1}{2} \sqrt{2 \abs{E_b}} < 2 \delta^{1/4} \sqrt{m}.$$
    Since $r\ge 3$ and $\lambda(G) \geq \frac{10}{9} \sqrt{m}$, we conclude that
    $$s_t \leq 0.9 s_{t - 1} + 2\delta^{1/4}.$$
    Together with the initial condition $s_1 \leq 1$, we can solve this recurrence and obtain that for every $t \leq \delta^{-1/8} / 2$, 
    $$s_t \leq 20 \delta^{1/4} + 0.9^{t - 1}.$$
    Taking $\ell := \floor{\delta^{-1/8}/2}$, since  $\delta$ is sufficiently small, we get $0.9^{\ell -1} < 10 \delta^{1/4}$. Thus, 
    \begin{equation} \label{eq-sl}
    s_{\ell} \leq 30 \delta^{1/4}.
    \end{equation}
    
    Let $V':=V(G)\setminus \Gamma_{\ell}$ and  $G':=G[V']$. 
    In what follows, we verify that $G'$ has the desired properties. 
    Firstly, by the definition of $V'$, we know that for every $i \in V'$,  
    $$x_i \geq (1 - 2\delta^{1/4} \ell) \alpha \geq (1 - 2\delta^{1/8}) \alpha.$$
    Recall that $\bm{x}$ is a unit vector and 
    $\alpha = (\frac{r-1}{r}\cdot \frac{1}{2m})^{1/4}$. 
    Thus, we get 
    $$
    |V'|   \leq \frac{1}{(1 - 2\delta^{1/8})^2 \alpha^2} 
    \sum_{i\in V'} x_i^2 \leq \frac{1}{(1 - 2\delta^{1/8})^2 \alpha^2} \leq (1 + \varepsilon) \cdot \sqrt{\frac{2rm}{r - 1}}.
$$    
Secondly, let $R$ be the set of edges with at least one vertex in $\Gamma_{\ell}$. Recall by~\eqref{eq-expect} that
    $$\sum_{ij \in R} \left(x_i x_j - \alpha^2\right)^2 \leq \sum_{ij \in E} \left(x_i x_j - \alpha^2\right)^2 \leq \delta.$$
    By the Cauchy--Schwarz inequality, we have
    $$ |R| \alpha^2 - \sum_{ij\in R} x_ix_j = 
    \sum_{ij \in R} (\alpha^2 - x_i x_j) \leq \sqrt{\delta \abs{R}}.$$
    On the other hand, we have
    $$\sum_{ij \in R} x_i x_j \leq \sum_{i\in \Gamma_{\ell}} x_i \sum_{j\in N(i)} x_j = 
    \lambda(G) \sum_{i \in \Gamma_{\ell}} x_i^2 < 30 \sqrt{2m} \delta^{1/4},$$
    where the last inequality holds by (\ref{eq-sl}). 
    So we conclude that
    $$\frac{1}{2\sqrt{m}}\abs{R} \leq \alpha^2 \abs{R} < 30 \sqrt{2m} \delta^{1/4} + \sqrt{\delta \abs{R}}.$$
    Solving this inequality, we obtain  $\abs{R} \leq 200\delta^{1/4} m < \varepsilon m$, which completes the proof.
\end{proof}

Now, we are ready to prove \Cref{thm-stability-clique-4} 
by combining Lemmas \ref{lem:nikiforov-stability-clique} with \ref{thm-niki-stability}. 

\begin{proof}[{\bf Proof of \Cref{thm-stability-clique-4}}] 
Let $\varepsilon>0$. Take $\varepsilon'>0$ sufficiently small so that $\varepsilon'<\varepsilon/2$ and $3(\varepsilon')^{1/4}<\delta_{\ref{thm-niki-stability}}(\varepsilon/8)$. Set $\delta=\delta_{\ref{lem:nikiforov-stability-clique}}(\varepsilon')$. Assume that $G$ is an $m$-edge $K_{r+1}$-free graph with $\lambda^2 (G)\ge {(1- \frac{1}{r} - \delta)2m}$. 
Applying Lemma \ref{lem:nikiforov-stability-clique} to $G$ with $\varepsilon_{\ref{lem:nikiforov-stability-clique}}=\varepsilon'$, we obtain a subgraph $G'$ on vertex set $V'$ with $|V'|\le (1+\varepsilon')\sqrt{2m/(1-1/r)}$ and $d(G,G')\le \varepsilon' m$.
Then we can apply the vertex-spectral stability, \Cref{thm-niki-stability}, to the $K_{r+1}$-free graph $G'$. 
Indeed, using the Rayleigh's formula, we have 
\[ \lambda (G')\ge \sqrt{\left(1- \frac{1}{r} - \delta \right)2m} - \sqrt{2\varepsilon' m} 
\ge \sqrt{\left(1- \frac{1}{r} - 3\sqrt{\varepsilon'} \right)2m}. \]
Then using the upper bound on $|V'|$, we can compute that 
\[ \lambda (G') \ge \left( 1-\frac{1}{r} - 3(\varepsilon')^{1/4}\right) |V'|. \]

We can choose $m$ sufficiently large such that $|V'|> \lambda (G') \ge \sqrt{m} \ge (n_0)_{\ref{thm-niki-stability}}$. 
Therefore, applying \Cref{thm-niki-stability} with $\varepsilon_{\ref{thm-niki-stability}}=\varepsilon/8$, 
we know that $G'$ differs from the $r$-partite Tur\'{a}n graph 
$T_{V',r}$ in at most 
$\varepsilon|V'|^2/8\le \varepsilon m/2$ edges. 
So it follows that 
$$d(G,T_{V',r}) \le d(G,G') + d(G',T_{V',r}) \le \varepsilon'm+\varepsilon m/2\le \varepsilon m,$$
as desired.
\end{proof}




%=========

\subsection{General graphs: Proving \Cref{thm:nikiforov-stability}}\label{sec:F-general}
In this subsection, we shall employ our framework in the proof of \Cref{conj-LLF} to complete the proof of \Cref{thm:nikiforov-stability} for general graphs $F$. 
First, we handle the case when $\chi (F)=r+1$ with $r \geq 3$. The main step is to find a large subgraph $G'$ retaining most of the edges such that its Perron--Frobenius eigenvector satisfies $x_i'x_j'\ge \delta m^{-1/2}$ for each $ij\in E(G')$. This in conjunction with Lemma \ref{lem:max} implies that $G'$ is dense: it has $O(\sqrt{m})$ vertices. We can then apply Lemma \ref{lemEFR} to $G'$ to obtain a $K_{r+1}$-free subgraph $G''$ by removing $o(|G'|^2)$ edges. This would complete the proof as $\lambda (G'')=(1-\frac{1}{r} -o(1))2m$ and we reduce the problem to clique case when $F = K_{r + 1}$.

We now implement this strategy in detail. To find the desired $G'$ mentioned above, we need the following lemma about the stability of the spectral radius after removing few edges. It eliminates some ``bad'' edges of $G$ that have small contribution to $\lambda(G)$. 


\begin{lemma}
\label{lem:removing-bad-edges}
Let $F$ be a graph with $\chi (F)=r+1\ge 3$. For every $\varepsilon \in (0, 0.01)$, there exist $m_0$ depending only on $F$ and $\varepsilon$ such that the following holds. If $G$ is an $F$-free graph with $m\ge m_0$ edges and $\lambda^2 (G)\ge (1- \frac{1}{r} - \varepsilon) 2m$, then $G$ contains a subgraph $G'$ with  $m'\ge (1 - 4\eps) m$ edges such that the unit Perron--Frobenius eigenvector $\bm{x}'$ of $G'$ satisfies that for every  $ij \in E(G')$, 
$$x'_i x'_j \geq \varepsilon m^{-1/2}.$$
\end{lemma}
\begin{proof}
    Let $\bm{x}=(x_i)_{i\in V}$ be the unit Perron--Frobenius eigenvector of $G$. 
    Suppose that $e = ij$ is an edge of $G$, then the Rayleigh quotient implies 
    $$\lambda(G \backslash e) \geq \lambda(G) - 2x_i x_j.$$ 
    We start with $G' = G$, and repeatedly delete edges $ij$ in $G'$ with $x'_i x'_j < \varepsilon m^{-1/2}$ from $G'$. Suppose on the contrary that we can repeat this deletion for $\ceil{4\varepsilon m}$ steps. Then the resulting graph $G'$ has $m' \leq (1 - 4\varepsilon)m$ edges, and satisfies
    $$\lambda(G') \geq \lambda(G) - 2\varepsilon m^{-1/2} \cdot \ceil{4\varepsilon m} \geq \lambda(G) - 9\varepsilon^2 m^{1/2}.$$
    Squaring the above inequality, we have 
    $$\lambda^2(G') > \lambda^2(G) - 18\varepsilon^2 m^{1/2} \lambda(G) > 
    \left(1- \frac{1}{r} - \varepsilon \right) 2m - 36 \varepsilon^2 m > 
    \left(1- \frac{1}{r} + \frac{\varepsilon}{4} \right) 2m'.$$
    However, $G'$ is $F$-free. For sufficiently large $m$, this contradicts with \Cref{conj-LLF}.
\end{proof} 

Now, we are ready to prove the edge-spectral stability in  \Cref{thm:nikiforov-stability} when $r\ge 3$. 

\begin{proof}[{\bf Proof of part (b) of \Cref{thm:nikiforov-stability}}]
    Let $\delta_0$ be the value of $\delta$ for \Cref{thm-stability-clique-4} with $\varepsilon' = \varepsilon / 3$. Without loss of generality, we may assume that $\delta_0 < \min\{10^{-10},  \varepsilon/10\}$. We take $\delta = \delta_0^3$. Let $G$ be an $F$-free graph with $m\ge m_0$ edges and $\lambda^2 (G)\ge (1- \frac{1}{r} - \delta) 2m$.
    
    We shall throw away few edges of $G$ to yield a large $K_{r+1}$-free subgraph. Applying \Cref{lem:removing-bad-edges} to $G$, we find a subgraph $G'$ of $G$ with  $m'\ge (1 - 4\delta) m$ edges such that the Perron--Frobenius eigenvector $\bm{x}'=(x_i')_{i\in V(G')}$  satisfies $x'_i x'_j \geq \delta m^{-1/2}$ for every edge $ij$ of $G'$. Removing all the isolated vertices from $G'$, let $V'$ be the vertex set of $G'$. By \Cref{lem:remove-edges}, we have
    $$\lambda^2(G') > \lambda^2(G) - 4\sqrt{4\delta} m > \left(1- \frac{1}{r} - 5\sqrt{\delta} \right) 2m.$$
    Since $r \geq 3$, we have $\lambda^2(G')  \geq 1.3 m$. By Lemma \ref{lem:max}, we have 
    \[  \max\{x'_i:i \in V(G') \} \leq 124 (m')^{-1/4} \leq 125 m^{-1/4}. \] 
     Thus, we have 
     \[ \min\{x'_i :i \in V(G')\} \ge 125^{-1}\delta m^{-1/4} . \] 
     Note that $1=\sum_{i\in V'} x_i'^2 
     \ge |V'| 125^{-2}\delta^2m^{-1/2}$, 
     which implies $\abs{V'} \le 125^2 \delta^{-2} m^{1/2}$.

    Applying Lemma \ref{lemEFR}, 
    there exists an $m_0$ such that 
    for $m \ge m_0$, we can remove at most $125^{-4}\delta^5 \cdot |V'|^2 \le \delta m$ edges from $G'$ to form a new subgraph $G''$ that is $K_{r + 1}$-free. By \Cref{lem:remove-edges}, we have
    $$\lambda^2(G'') > \lambda^2(G) - 4\sqrt{5\delta} m > \left(1- \frac{1}{r} - 5 \sqrt{\delta} \right) 2m > \left(1- \frac{1}{r} - \delta_0 \right) 2m.$$
    Applying \Cref{thm-stability-clique-4}, there exists a vertex set $C \subseteq V(G'')$ such that
    $$d(G'', T_{C, r}) \leq \varepsilon m / 3.$$
    Thus we obtain
    \begin{align*}
    d(G, T_{C, r}) &\leq  d(G, G')+ d(G',G'')  +d(G'', T_{C, r}) \\
    &\leq 4 \delta m + \delta m +  \varepsilon m / 3  \\
    & \leq \varepsilon m , 
    \end{align*}
    as desired.
\end{proof}



For the case $r = 2$,  we cannot apply Lemma \ref{lem:max} to conclude that all the coordinates $(x_v)_{v\in V}$ are bounded by $O(m^{-1/4})$. To overcome this obstacle, we use instead the following lemma.

\begin{lemma} \label{lem:max-2}
For every $\varepsilon > 0$, there exists $\delta > 0$ such that the following holds. Let $G$ be a graph with $m$ edges, and let $\bm{x}$ be the unit Perron--Frobenius eigenvector of $G$. If $\lambda^2 (G) \ge (1 - \delta)m$ and 
\begin{equation} \label{eq-xv-large}
\max\{x_i : i\in V(G)\} > {\delta}^{-1} m^{-1/4}, 
\end{equation}
then there are disjoint vertex subsets $U, V$ of $G$ such that $d(G, K_{U, V}) \leq \varepsilon m$.
\end{lemma}

\begin{proof}
Using (\ref{eq-xv-large}), we shall show that $G$ contains a bipartite subgraph $G'$ with large spectral radius. Then we can invoke  \Cref{lem:close-to-Kab} to complete the proof. 
We set $\delta := (\delta_0/5)^8$,  
where $\delta_0=(\varepsilon/2)^4/100$ is the constant from \Cref{lem:close-to-Kab}. 
Let $\lambda = \lambda (G)$ and 
$x_i= \max\{x_v : v\in V(G)\}$. 
Since $\lambda x_i = \sum_{j\in N(i)} x_j$, we have 
\begin{equation}\label{eq-eigen}
  \lambda^2 x_i = \sum_{j\in N(i)} \lambda x_j= \sum_{j\in N(i)} 
\sum_{k\in N(j)} x_k \le \sum_{jk \in E(G)}(x_j +x_k). 
\end{equation}
Let 
$$R := \{ j\in V(G): x_j > \delta^{-1/2}m^{-1/4}\}.$$ 
Since 
$\sum_{v\in V} x_v^2 =1$, we see that 
\[ \abs{R} <  \delta m^{1/2}. \] 
We obtain from (\ref{eq-eigen}) that 
$$ (1-\delta)mx_i\le \lambda^2 x_i \leq \sum_{jk \in E(R)}(x_j +x_k) +  \sum_{ jk \in E(\bar{R})}(x_j +x_k) +  \sum_{jk \in E(R, \bar{R})}(x_j +x_k).$$
The first term is at most $\abs{R}^2 x_i \leq \delta^2 m x_i$. By the definition of $R$ and that $x_i > \delta^{-1}m^{-1/4}$ due to \eqref{eq-xv-large}, the second term above is at most $m \cdot 2\delta^{-1/2}m^{-1/4} \leq 2\delta^{1/2} m x_i$. Thus we have
$$\sum_{jk \in E(R, \bar{R})}(x_j +x_k) \geq (1 - 4\delta^{1/2}) m x_i.$$
For each edge $jk \in E(R, \bar{R})$, we have $x_j + x_k \leq x_i + \delta^{-1/2} m^{-1/4} \le 
(1 + \delta^{1/2})x_i$. Then 
\[ e(R,\bar{R})\cdot  (1+ \delta^{1/2})x_i 
\ge (1- 4\delta^{1/2} ) mx_i. \]
We conclude that at least $(1 - 5\delta^{1/2})m$ edges lies between $R$ and $\bar{R}$. Now, we denote $G'=G[R,\bar{R}]$. 
Then \Cref{lem:close-to-Kab} can be invoked on $G'$ to complete the proof. 
Indeed, Rayleigh's formula yields 
\[ \lambda (G') \ge \lambda (G) - 
\lambda (G\setminus G') 
\ge \sqrt{(1-\delta) m} - \sqrt{10\delta^{1/2}m} 
> (1- 5\delta^{1/8}) \sqrt{m}. \]
Note that $G'$ is bipartite. 
By \Cref{lem:close-to-Kab}, 
there exist two disjoint vertex subsets $U,V$ such that $d(G',K_{U,V})\le \varepsilon m/2$. 
Thus, we have 
\[ d(G,K_{U,V}) \le d(G,G') + d(G',K_{U,V}) \le 5\delta^{1/2}m + \varepsilon m/2 \le \varepsilon m. \]
This completes the proof.
\end{proof}

We now finish the proof of \Cref{thm:nikiforov-stability} 
for the case $\chi (F)=r+1=3$. 


\begin{proof}[{\bf Proof of part (a) of \Cref{thm:nikiforov-stability}}]
    Let $\delta_0$ be the value of $\delta$ from \Cref{thm-stability-triangle} with $\varepsilon' = \varepsilon / 3$, and let $\delta_1$ be the value of $\delta$ from \Cref{lem:max-2} with $\varepsilon' = \varepsilon / 2$. Without loss of generality, we may assume that $\delta_0 < \min\{10^{-10}, \varepsilon/10 \}$. We take $\delta = \min\{\delta_0^3, \delta_1^3\}$. Let $G$ be an $F$-free graph with $m\ge m_0$ edges and $\lambda^2 (G)\ge (1 - 2\delta) m$. 

    Applying \Cref{lem:removing-bad-edges} to $G$, we find a subgraph $G'$ of $G$ with at least $(1 - 4 \delta) m$ edges such that the Perron--Frobenius eigenvector $\bm{x}'=(x_i')_{i\in V(G')}$ satisfies $x'_i x'_j \geq \delta m^{-1/2}$ for every  $ij\in E(G')$. Removing all isolated vertices from $G'$, let $V'$ be the vertex set of $G'$. By \Cref{lem:remove-edges}, we have
    $$\lambda^2(G') > \lambda^2(G) - 4\sqrt{4\delta} m > (1- 10\sqrt{\delta}) m.$$
    We now divide the remaining proof into two cases.

\smallskip
    
    \noindent\textbf{Case 1.} $\max\{x_v' : v\in V(G')\} > \delta_1^{-1} m^{-1/4} $. 
    
    Using \Cref{lem:max-2}, there exist disjoint vertex subsets $U, V$ such that      
    $d(G',K_{U, V}) \le \varepsilon m/2$. 
    Thus, it follows that 
    $d(G,K_{U, V}) \le d(G,G') + d(G' ,K_{U, V}) \le 4 \delta m + \varepsilon m/2 \le \varepsilon m$ as desired. 

\smallskip
    
    \noindent\textbf{Case 2.} $\max\{x_v' : v\in V(G')\} \leq \delta_1^{-1} m^{-1/4}$. 
    
    In this case, we can apply a similar argument as in the case $r \geq 3$. 
    First of all, we have 
    $\min\{x_v': v\in V(G')\} \ge \delta \delta_1 m^{-1/4}$. 
    Observe that $1=\sum_{i\in V'} (x_i')^2\ge |V'| \cdot (\delta \delta_1)^2 m^{-1/2}$, which gives 
     $|V'| \le (\delta \delta_1)^{-2} m^{1/2} $. 
    By Lemma \ref{lemEFR}, there exists $m_0$ such that for $m\ge m_0$, 
    we can remove at most 
    $(\delta \delta_1)^5|V'|^2 \le 
    \delta m$ edges from $G'$ to get a 
    triangle-free graph $G''$. By 
    \Cref{lem:remove-edges}, we get 
    \[  \lambda^2(G'') > \lambda^2(G) - 
    4\sqrt{5\delta} m > (1- 11\sqrt{\delta})m. \]
    Applying \Cref{thm-stability-triangle}, there exist disjoint vertex subsets $U, V$ such that      
    $d(G',K_{U, V}) \le \varepsilon m/3$. Thus, we get  
    $ d(G,K_{U, V}) \le d(G,G') + d(G',G'') + 
    d(G'',K_{U, V}) \le \varepsilon m$, as needed. 
\end{proof}







\section{Maximizing common neighbors}

\label{sec:common}

We first give a random construction to show that our bound on $t$ in~\Cref{thm-common} is the best possible. Let $\varepsilon>0$ and consider the binomial random graph $G(n, \frac{1}{2} + \varepsilon)$. 
Standard concentration bounds show that for sufficiently large $n$, with high probability both of the following events occur:
\begin{enumerate}
    \item $G$ has at least $(\frac{1}{2} + \frac{\varepsilon}{2}) \binom{n}{2}$ edges;
    \item the codegree of any two vertices is at most $(\frac{1}{2} + 2\varepsilon)^2 n$.
\end{enumerate}
Consider any $G$ such that both the above two properties hold. Then first one implies that $m > n^2 / 4$, and so $\lambda(G) \geq \frac{2m}{n} > \sqrt{m}$. The second implies that if $K_{2, t}$ is a subgraph of $G$, then $t \leq (\frac{1}{2} + 2\varepsilon)^2 n \leq \left(\frac{1}{2} + 6\varepsilon\right) \sqrt{m}$. This random construction reveals the tightness of \Cref{thm-common}. 

\medskip 
To prove~\Cref{thm-common}, we need the following lemma.  

\begin{lemma} \label{lem-lower-bound}
    Let $G=(V,E)$ be a graph with $m$ edges and 
    $\bm{x}$ be the unit Perron--Frobenius eigenvector of $G$. Then 
    $$\sum_{i,j\in V} |N(i)\cap N(j)|\cdot x_i^2 x_j^2 \geq \frac{\lambda^3(G)}{2m}.$$
\end{lemma}

The bound in Lemma \ref{lem-lower-bound} 
is the best possible when we take $G=K_{\sqrt{m},\sqrt{m}}$. 

\begin{proof}
Let $\lambda =\lambda (G)$ and $C_{i, j}$ be the number of common neighbors of $i$ and $j$. In particular, we have that  $C_{i, i} $ is the degree $ d_i$. 
We can write
$$\sum_{i, j\in V} C_{i, j} x_i^2 x_j^2 
= \sum_{i\in V} x_i^2 \sum_{k \in N(i)} \sum_{j \in N(k)} x_j^2.$$
By the Cauchy--Schwarz inequality, we have
$$\sum_{j \in N(k)} x_j^2 \geq \frac{1}{d_k} \left(\sum_{j \in N(k)} x_j\right)^2 = \frac{\lambda^2}{d_k} x_k^2.$$
Thus, we obtain 
$$\sum_{i, j\in V} C_{i, j} x_i^2 x_j^2 \geq 
\lambda^2 \sum_{i\in V} \sum_{k \in N(i)} \frac{x_i^2 x_k^2}{d_k} = 2\lambda^2 \sum_{ik \in E} \frac{x_i^2 x_k^2}{d_k}.$$
By symmetry, we have 
\[ \sum_{i, j\in V} C_{i, j} x_i^2 x_j^2 \geq 2\lambda^2 \sum_{ik \in E} \frac{x_i^2 x_k^2}{d_i}.\] 
Adding the above two inequalities and using the AM-GM inequality, we get 
\begin{equation} \label{eq-low}
    \sum_{i, j\in V} C_{i, j} x_i^2 x_j^2 \geq 2\lambda^2  \sum_{ik \in E} \frac{x_i^2 x_k^2}{\sqrt{d_i d_k}}.
\end{equation}
By the Cauchy--Schwarz inequality again, we have
\begin{equation*} 
    \sum_{ik \in E} \frac{x_i^2 x_k^2}{\sqrt{d_i d_k}} \geq \frac{\left(\sum_{i k \in E} x_i x_k\right)^2}{\sum_{ik \in E} \sqrt{d_i d_k}}.
\end{equation*}
On the one hand, it follows from the definition that $\sum_{ik \in E} x_i x_k = \lambda / 2.$
On the other hand, the vector $ (\sqrt{d_i})_{i=1}^n$ has length $\sqrt{2m}$. Therefore, the Rayleigh formula gives 
$$\sum_{ik \in E} \sqrt{d_i d_k} \leq 
\frac{\lambda}{2} \sum_{i=1}^n d_i =
\frac{\lambda}{2} \cdot 2m .$$
Thus, one more application of the Cauchy--Schwarz inequality yields
\begin{equation} \label{eq-bow}
   \sum_{ik \in E} \frac{x_i^2 x_k^2}{\sqrt{d_i d_k}} \geq \frac{\lambda}{4m}. 
\end{equation}
Combining  (\ref{eq-low}) 
with (\ref{eq-bow}), we conclude that $\sum_{i, j\in V} C_{i, j} x_i^2 x_j^2 \geq \frac{\lambda^3}{2m}$ as desired.
\end{proof}



Now, we are in a position to prove Theorem \ref{thm-common}. 

\subsection{Proof of Theorem \ref{thm-common}}


Let $C_{i,j}=|N(i)\cap N(j)|$. Since $\lambda (G)> \sqrt{m}$, \Cref{lem-lower-bound} gives 
    $$\sum_{i\in V} d_i x_i^4 + \sum_{i \neq j} C_{i, j} x_i^2 x_j^2 > \frac{\sqrt{m}}{2}.$$
    We denote $t := \max_{i \neq j} C_{i, j}$. Then 
    $$\sum_{i \neq j} C_{i, j} x_i^2 x_j^2 \leq t \sum_{i \neq j} x_i^2 x_j^2 \leq t.$$
    Note that 
    $$\sum_{i\in V} d_i x_i^4 \leq 2m \cdot (\max_{i\in V} x_i)^4$$
    Therefore, we conclude that 
    \begin{equation} \label{eq-t-max}
        t + 2m \cdot (\max_{i\in V} x_i)^4 > \frac{\sqrt{m}}{2}.
    \end{equation}
    
    Let $\varepsilon >0$ be sufficiently small and 
    $\delta $ be the number determined in 
    the spectral stability Lemma \ref{lem:max-2}. 
    If $\max_{i\in V} \{x_i\} \le \delta^{-1}m^{-1/4}$, 
    then (\ref{eq-t-max}) implies $t> \sqrt{m}/2 - 2\delta^{-4}$, as desired. 
    
    We then assume $\max_{i\in V} \{x_i\} > \delta^{-1} m^{-1/4} $, so that Lemma \ref{lem:max-2} is applicable. 
    There exist disjoint vertex subsets $A,B$ of $G$ such that 
    $d(G,K_{A,B})\le \varepsilon m$, which means that the total number of edges with two endpoints within $A$ or $B$, the missing edges between $A$ and $B$, and the edges with at least one endpoint outside $A\cup B$ are at most $\varepsilon m$. 
    Then $e(A)+e(B)\le \varepsilon m$ and 
    \[ (1- \varepsilon)m \le  |A||B|  \le (1+ \varepsilon )m. \]  
    We may assume that $|A|\le |B|$. 
    Then $|A| \le \sqrt{(1+\varepsilon)m}$. 

We claim that $ |A|\ge 2$. Suppose on the contrary that 
$|A|=1$. Then using Lemma \ref{lem-Delta}, we know that 
$ e(G)\le \Delta (G) + \varepsilon m \le m/2 + m^{0.99} + \varepsilon m <m$, which is a contradiction. 
 Thus, we must have $|A|\ge 2$. 
By double counting, we can show that 
    there are two vertices in $A$ that have at least ${(1- 10\varepsilon)m}/{|A|}$ common neighbors in $B$. 
Indeed, suppose on the contrary that 
for any $2$-set $\{i,j\} \subseteq A$, we have $C_{i,j}< {(1-10\varepsilon)m}/{|A|}$. Then 
\[ \sum_{\{i,j\} \subseteq A} C_{i,j} < {|A| \choose 2} \cdot \frac{(1-10\varepsilon)m}{|A|} .\]
On the other hand, the Jensen inequality gives  
\[ \sum_{\{i,j\}\subseteq A} C_{i,j} = \sum_{k\in B} {d_k \choose 2} \ge |B| {|B|^{-1} \sum_{k\in B}d_k \choose 2} .  \]
Combining these two bounds, we get 
${|A| \choose 2} \frac{(1-10\varepsilon)m}{|A|} > |B| {m/|B| \choose 2}$, which yields 
$(|A| -1) (1- 10\varepsilon) > \frac{m}{|B|}-1$. 
Since $|A| \ge 2$ and $|A||B| \le (1+\varepsilon)m$, 
by simplifying, we have $m-10(1+\varepsilon)m + 10 |B|>0$, which leads to a contradiction since $0<\varepsilon< 0.01$ and $|B|\le \frac{1+\varepsilon}{2}m$.  

Since $|A|\le \sqrt{(1+\varepsilon)m}$, we have $t\ge (1- 10\varepsilon)m/|A| \ge (1-5\sqrt{\varepsilon })\sqrt{m}$, as desired. This completes the proof.

\section{Concluding remarks}\label{sec:concluding}
In this paper, we obtained an edge-spectral extension of Erd\H{o}s--Stone--Simonovits theorem and its stability result. As an application of our edge-spectral stability result, we proved an optimal bound on the size of $K_{2,t}$ in $m$-edge graphs with spectral radius larger than $\sqrt{m}$. For graphs with spectral radius larger than $(1 + \varepsilon)\sqrt{m}$, our technique also finds, for any fixed $s$, a copy of $K_{s,t}$ with $t=\Omega(\sqrt{m})$.

\begin{theorem} \label{thm-squar-root}
For every $s\ge 2$ and $\varepsilon >0$, 
there exists $\delta = \delta (s,\varepsilon)>0$ such that 
if $G$ is an $m$-edge graph with 
$\lambda (G)\ge (1+ \varepsilon )\sqrt{m}$, then $G$ contains a copy 
of $K_{s,t}$ with $t\ge \delta \sqrt{m}$. 
\end{theorem}

We need the following variant of K\H{o}v\'{a}ri--S\'{o}s--Tur\'an's theorem. 

\begin{theorem} \label{prop-linear-size}
If $\varepsilon >0$ and $G$ is an $n$-vertex graph with $e(G)\ge \varepsilon n^2,$
then $G$ has a copy of 
$K_{s,t}$ with $s\ge 1$ and $t= (\varepsilon/{2s})^s n$. 
\end{theorem} 

\begin{proof}[Proof of~\Cref{thm-squar-root}]
We denote $\delta := (8\varepsilon^{16}/s)^{s}\sqrt{2}$. 
Assume on the contrary that 
$G$ is a graph with $m$ edges such that 
$\lambda (G)\ge (1+ \varepsilon )\sqrt{m}$ and
$G$ does not contain a copy of $K_{s,t}$, where 
$t=\delta \sqrt{m}$. 
Among such graphs, we choose a graph $G$ such that $|E(G)| + |V(G)|$ is minimal. 


We claim that $x_ix_j \ge 
1 / (8\sqrt{m})$ for every edge $ij\in E(G)$. 
Suppose on the contrary that there exists an edge $ij\in E(G)$ satisfying $x_ix_j < 1 / (8\sqrt{m})$.  
Removing the edge $ij$ from 
$G$, we obtain a subgraph $G'$ such that 
\[ \lambda (G')\ge \lambda (G) - 2x_ix_j 
>(1+\varepsilon)\sqrt{m} - 1 / (4\sqrt{m}) 
> (1+\varepsilon)\sqrt{m-1}. \]  
Clearly, $G'$ contains no copy of $K_{s,t}$. 
So $G'$ contradicts with the minimality of $G$.


Note that $\lambda^2(G)\ge (1+2\varepsilon)m$. 
By Lemma \ref{lem:max}, we know that 
$\max\{x_v: v\in V(G)\} < 1/(16\varepsilon^4 m^{1/4})$. 
Therefore, we get $\min\{x_v: v\in V(G)\} 
> 2\varepsilon^4/m^{1/4}$. 
Then $1=\sum_{i\in V} x_i^2 > n \cdot 4\varepsilon^8/\sqrt{m}$. It follows that  
 $G$ is a dense graph with 
$m>16 \varepsilon^{16} n^2$ edges and a desired copy of $K_{s,t}$ is guaranteed by the K\H{o}v\'{a}ri--S\'{o}s--Tur\'an theorem,  contradicting with the assumption. 
\end{proof}

 
 A similar argument of \Cref{thm-squar-root} gives the following result.  
 
 \begin{theorem} \label{thm-log-m}
For every $\varepsilon >0$, 
there is a constant $\delta = \delta (\varepsilon) >0$ such that 
if $G$ is an $m$-edge graph with 
$\lambda (G)\ge (1+ \varepsilon )\sqrt{m}$, then $G$ contains a copy 
of $K_{t,t}$ with $t= \delta \log m$. 
\end{theorem} 

\begin{proof}
We denote by $\delta:=\left(2\log (1/32\varepsilon^{16}) \right)^{-1}$. 
Assume on the contrary that $G$ does not contain a copy of $K_{t,t}$ with $t=\delta \log m$. Similar to the proof of Theorem \ref{thm-squar-root}, we replace \Cref{prop-linear-size} with the following density result: 
If $G$ is an $n$-vertex graph with 
$ e(G)\ge \varepsilon {n \choose 2}$, where $\varepsilon >0$, 
 then $G$ contains a copy of $K_{t,t}$ where $t=\lfloor {\log n}\cdot \left({2 \log (1/ \varepsilon)} \right)^{-1} \rfloor$. Therefore, the density result 
 $m>32 \varepsilon^{16} {n \choose 2}$ implies that 
 $G$ contains 
    a copy of $K_{t,t}$ with $t\ge {\log n}\cdot (2\log (1/32\varepsilon^{16}) )^{-1}> 
    {\log (2m)} \cdot (2{\log (1/32\varepsilon^{16})} )^{-1}$. 
    This contradicts with the assumption. 
\end{proof}

The $\log m$ bound above is the best possible by considering the random graph $G(n, p)$, where each edge is presented independently with probability $p$. By a standard probabilistic argument, for any $p\in (0,1)$, there exists an $n$-vertex graph $G$ with $e(G)\ge p{n \choose 2}$, but $G$ does not contain $K_{t,t}$ with $t=\frac{3\log n}{\log (1/p)}$. Taking $p=\frac{2}{3}$, we have $\lambda (G) > 1.1 \sqrt{m}$, but $G$ has no $K_{t,t}$ with $t=8\log m$.  


\medskip 
There are many promising directions to explore in the edge-spectral Tur\'{a}n type problem, and we provide another example for cycles with consecutive lengths. 

\begin{theorem}
    For any $\varepsilon >0$, there exists $\delta = \delta (\varepsilon) >0$ such that if $G$ is an $m$-edge graph with $\lambda (G)\ge (1+ \varepsilon )\sqrt{m}$, then $G$ contains a cycle of length $\ell$ for every $3\le \ell \le \delta \sqrt{m}$. 
\end{theorem}


\begin{proof}
    We denote $\delta =2 \varepsilon$ and $\lambda = \lambda (G)$. 
    Then $\lambda (G)\ge (1+\varepsilon)\sqrt{m}$ gives 
    $\lambda^2 - \varepsilon \sqrt{m} 
    \lambda > m$. Now, we consider the matrix $T:=A^2 - \varepsilon \sqrt{m} A$. Since $A$ is non-negative, 
    it is well-known that 
    \[ \lambda (T)\le \max_{i\in V} \sum_{j\in V} t_{ij}.\]  Clearly, we have 
    $\lambda (T) = \lambda^2- \varepsilon \sqrt{m} \lambda > m$. 
    On the other hand, we get  
 \[ \sum_{j\in V}t_{ij} =
    \sum_{k\sim i} \sum_{j\sim k} 1 - \varepsilon \sqrt{m} \sum_{j\sim i} 1 
    \le m +e(N(i)) - \varepsilon \sqrt{m} |N(i)|. \]
    Therefore, there exists a vertex  
    $i\in V$ such that $e(N(i)) > \varepsilon \sqrt{m} |N(i)|$. 
    By the Erd\H{o}s--Gallai theorem, we find a path $P_{t+1}$ in $N(i)$ with length $t := 2 \lfloor \varepsilon \sqrt{m} \rfloor +1$. So $G$ contains a copy of $K_1\vee P_{t+1}$, which implies that $G$ contains all cycles $C_{\ell}$ with $3\le \ell \le 2\varepsilon \sqrt{m}$, as needed.     
\end{proof}

The above bound $\Omega (\sqrt{m})$ on $\ell$ is optimal up to a constant factor.




%\section*{Declaration of competing interest}
%The authors declare that they have no conflicts of interest to this work.


%\section*{Data availability}
%No data was used for the research described in the article.


\section*{Acknowledgements}
This paper was dedicated to Vladimir Nikiforov, whose beautiful works on spectral extremal graph theory inspire the authors. 
This work was initiated after the 1st IBS ECOPRO student research program in fall 2023. The first and third authors would like to thank ECOPRO group for hosting them. The authors also thank Prof. Hao Huang for motivating this project. 




\frenchspacing

%\bibliographystyle{abbrv}
%\bibliography{ref.bib}

\input{Part-1-spectral-stability.bbl}

\appendix



\section{An alternative proof of \Cref{thm-stability-clique-4}} 

\label{sec-A}


In this section, we provide an alternative proof of the edge-spectral stability result for cliques by using the Kruskal--Katona stability theorem in \cite{LM2023}. We provide a detailed argument when forbidding a $K_4$, which easily extends to larger cliques. Let us first outline the idea.

\begin{enumerate}
    \item If $\lambda(G)$ is close to the maximal, that is, 
    $\lambda (G)= (1-o(1))\sqrt{4m/3}$, then we can conclude that $k_3(G)$ is close to its maximal value $(m/3)^{3/2}$ in the S\'{o}s--Straus theorem.
    \item We construct a hypergraph $\mathcal{H}$ whose edges are the triangles of $G$. By Theorem 1.8 in \cite{LM2023}, 
    we see that $\mathcal{H}$ is close to a complete $r$-partite hypergraph.
    \item Returning to the $2$-graphs, 
    it is not hard to show that $G$ is close to $T_{n, r}$, as desired. 
\end{enumerate}

To begin with, we need some notations. 
We denote by $k_r(G)$ the number of cliques of order $r$ in $G$.  
For example, we have $k_2(G)=m$ and $k_1(G)=n$. 

\begin{lemma}[Nikiforov \cite{Niki2002}] \label{thm-Niki-2002-p}
Let $G$ be a $K_{\ell+1}$-free graph. Then 
\begin{equation} \label{eq-Niki-2002-clique}
 \lambda^{\ell} (G) \le 
\sum_{i=2}^{\ell} (i-1)k_i(G) \lambda (G)^{\ell -i},
\end{equation}
\end{lemma}

\begin{lemma}[S\'{o}s--Straus \cite{SS1982}; 
Fisher--Ryan \cite{FR1992}]   
\label{lem-SS-FR}
Let $G$ be a $K_{\ell+1}$-free graph on $n$ vertices. 
For every $i\in [\ell]$, let $k_i$ denote the number of copies of $K_i$ 
in $G$. Then 
\begin{equation} \label{eq-SS-FR}
  \left( \frac{k_{\ell}}{{\ell \choose \ell}}\right)^{1/\ell} 
\le  \left( \frac{k_{\ell-1}}{{\ell \choose \ell-1}}\right)^{1/(\ell-1)} 
\le \cdots \le \left( \frac{k_2}{{\ell \choose 2}}\right)^{1/2} 
\le \left( \frac{k_1}{{\ell \choose 1}}\right)^{1/1}. 
\end{equation} 
\end{lemma}

In fact, we can obtain an alternative proof of Nikiforov's bound (\ref{eq-Niki-2002-cpc}) by using Lemmas \ref{thm-Niki-2002-p} and \ref{lem-SS-FR}. 
To show the stability of Nikiforov's bound, we need to establish a stability 
result of Lemma \ref{lem-SS-FR} for our purpose. This can be deduced from a stability result on hypergraphs. 

\begin{defn}
Let $\mathcal{H}$ be an $r$-uniform hypergraph on vertex set $V$. 
The shadow of $\mathcal{H}$ is defined as 
$ \partial \mathcal{H} = \left\{A \in \tbinom{V}{r-1} : \text{there is $B\in E(\mathcal{H})$ such that $A\subset B$}  \right\}$. 
\end{defn}

\iffalse 
In 2021, Liu and Mubayi \cite{LM2021} 
proved the following inequalities. 

\begin{theorem}[Liu--Mubayi, 2021] \label{thm-LM2021}
Let $\ell \ge r \ge 2$ and $\mathcal{H}$ be a $\mathcal{K}_{\ell +1}^{(r)}$-free 
$r$-graph. Then 
\[  \left( \frac{|\mathcal{H}|}{{\ell \choose r}} \right)^{\frac{1}{r}} \le 
\left( \frac{| \partial_1 \mathcal{H}|}{{\ell \choose r-1}} \right)^{\frac{1}{r-1}} . \]
\end{theorem}

This result can deduce the inequality in 
Lemma \ref{lem-SS-FR} by translating each $r$-clique of a $2$-graph into an edge of an $r$-graph.  
\fi 


In 2023, Liu and Mukherjee \cite{LM2023} proved 
the following stability result.  


\begin{theorem}[Liu--Mukherjee \cite{LM2023}] \label{thm-LM2022}
For every $\ell \ge r \ge 2$ and $\varepsilon >0$, 
there exist $n_0$ and $\delta >0$ such that 
the following holds for all numbers $x\ge n_0$. 
Suppose that $\mathcal{H}$ is an  $\mathcal{K}_{\ell +1}^{(r)}$-free $r$-graph with 
\[ |\partial \mathcal{H}| = {\ell \choose r-1} \left( \frac{x}{\ell}\right)^{r-1}~~\text{and}~~ 
|\mathcal{H}| \ge (1-\delta ){\ell \choose r} 
\left( \frac{x}{\ell}\right)^r. \]
Then there exists a vertex set $C\subseteq V(\mathcal{H})$ of size at most $\lceil x\rceil$
such that $d(\mathcal{H},\mathcal{T}_{C,\ell}^{(r)}) \le \varepsilon x^r$.  
\end{theorem}


Now, we are ready to show  
 \Cref{thm-stability-clique-4} for the case $r=3$. 

\begin{lemma}
    For any $0< \beta<0.01$, there exists $\delta >0$ such that if 
   $G$ is a $K_4$-free graph with $m$ edges and 
    $\lambda (G) \ge (1- \delta )(4m/3)^{1/2}$, then 
    $k_3(G)\ge (1- \beta)(m/3)^{3/2}$. 
\end{lemma}

\begin{proof}
We take $\delta = \beta^{10}$. 
Suppose on the contrary that  
\[ k_3(G)< (1- \beta) (m/3)^{3/2}. \]
Then Lemma \ref{thm-Niki-2002-p} gives 
\[  \lambda^3 (G) \le m\lambda(G) + 2k_3(G) < 
m \lambda (G) + 2 (m/3)^{3/2} - 2\beta (m/3)^{3/2}. \]
Next, we deduce a contradiction by showing that $\lambda(G) < (1- \delta) (4m/3)^{1/2}:=x_0$. We denote 
\[ f(x):=x^3 - mx - 2(m/3)^{3/2} + 2\beta (m/3)^{3/2}.  \]
Then we have $f(\lambda )<0$. To show $\lambda < x_0$, 
it suffices to prove that $f(x)\ge 0$ for every $x\ge x_0$. Firstly, we can observe that $(4m/3)^{1/2}$ is the largest root of 
$x^3 - m x - 2(m/3)^{3/2}=0$. Then 
\begin{align*}
    f(x_0) &= 
(1- \beta^{10})^3 (4m/3)^{3/2} - m (1- \beta^{10}) (4m/3)^{1/2} - 2(m/3)^{3/2} 
+ 2\beta (m/3)^{3/2}\\ 
&= (-3\beta^{10} + 3\beta^{20} - \beta^{30}) (4m/3)^{3/2} + \beta^{10} m (4m/3)^{1/2} +2\beta (m/3)^{3/2} \\
&= m^{3/2} \left(2\beta (1/3)^{3/2} + \beta^{10}(1/3)^{1/2} -(3\beta^{10} - 3\beta^{20}+ \beta^{30})(4/3)^{3/2} \right) \\
&>0.
\end{align*}
Secondly, we have $f'(x)=3x^2 - m$. Setting $f'(x)=0$, we get $x=\pm (m/3)^{1/2}$. Consequently, $f(x)$ is monotonically increasing for $x\ge (m/3)^{1/2}$. Thus, 
$f(x)\ge 0$ for every $x\ge x_0$. 
\end{proof}


The case of $K_4$-free graph in \Cref{thm-stability-clique-4} can be reduced as follows. 

\begin{lemma}
    For any $0<\varepsilon <0.01$, there is $\beta >0$ such that if $G$ is a $K_4$-free graph with $m$ edges and $k_3(G)\ge (1- \beta)(m/3)^{3/2}$, then $d(G,T_{C,3})\le \varepsilon  m$ for some vertex set $C\subseteq V(G)$. 
\end{lemma}

\begin{proof}
We denote $\gamma = \varepsilon /2$ and 
$\beta = \min\{\delta_0, \gamma^4\}$, where $\delta_0=\delta_0(\gamma)$ is defined as in \Cref{thm-LM2022}. 
    We construct a $3$-uniform hypergraph $\mathcal{H}$ in which $V(\mathcal{H})=V(G)$ and 
    $\{a,b,c\}$ is an edge of $\mathcal{H}$ if and only if $abc$ forms a triangle in $G$. 
    Then $|\mathcal{H}|=k_3(G)$ and $|\partial \mathcal{H}| \le e(G)=m$. 
    Since $G$ is $K_4$-free, we know by definition that $\mathcal{H}$ is $\mathcal{K}_4^{(3)}$-free. So $k_3(G)\ge (1- \beta )(m/3)^{3/2}$ gives 
    \[  |\mathcal{H}|^{1/3} \ge (1- \delta_0) 
    \left( {|\partial \mathcal{H}|}/{3}\right)^{1/2}. \]
    Setting $x:=3(\partial \mathcal{H}/3)^{1/2}$ in \Cref{thm-LM2022}, there exists 
     a vertex set $C\subseteq V(\mathcal{H})$ with 
    $|C|\le \lceil 3(m/3)^{1/2}\rceil$ such that 
    $\mathcal{H}$ is a subgraph of $\mathcal{T}_{C,3}^{(3)}$ after removing at most $\gamma^2 (m/3)^{3/2}$ edges. 

    
    We denote $C=V_1\sqcup V_2\sqcup V_3$. 
    Our goal is to show that such a partition is the desired $T_{C,3}$ for $G$. 
    We call a triangle of $G$ good if it intersects each $V_i$ in exactly one vertex. Otherwise, we call it bad. In other words, \Cref{thm-LM2022} implies that 
    $G$ has at most $\gamma^2(m/3)^{3/2}$ bad triangles. 
    
   Let $E_b$ be the set of edges in $G$ that are not contained in any $G[V_i,V_j]$ for $i\neq j$. We claim that $|E_b|\le \gamma m$. 
   Suppose on the contrary that $|E_b|> \gamma m$. Then 
    there are less than $(1-\gamma)m$ cross-edges in $V_1\sqcup V_2\sqcup V_3$. Since $G$ is $K_4$-free, by Lemma \ref{lem-SS-FR}, there are less than $((1-\gamma)m/3)^{3/2}$ good triangles in $G$. Therefore, it follows that 
    \[ k_3(G)<\gamma^2(m/3)^{3/2} + ((1-\gamma)m/3)^{3/2}. \]  
    By the assumption, we have 
    \[  k_3(G)\ge (1- \beta )(m/3)^{3/2} \ge (1- \gamma^4)(m/3)^{3/2}, \] 
     which leads to a contradiction since $1-\gamma^4 > \gamma^2 + (1-\gamma)^{3/2}$ for any $0<\gamma<0.74$. Thus, we must have $|E_b| \le \gamma m$, i.e.,  
    there are at most $\gamma m$ edges of $G$ outside $\bigcup_{i<j} G[V_i,V_j]$. 

   In what follows, we show that there are few missing cross-edges in the partition $C=V_1\sqcup V_2\sqcup V_3$. From the above argument, we know that $G$ has at least $(1- \beta -\gamma^2)(m/3)^{3/2}$ good triangles. 
    Let $E_s$ be the set of edges of $T_{C,3}$ that are missed in $G$.   
     Observe that each missing edge between $V_1$ and $V_2$ destroys $|V_3|$ triangles of $T_{C,3}$. 
    Since $\big| |V_i| - |V_j| \big|\le 1$, we have 
    \[ |C|/3 -1 \le |V_i|\le |C|/3 +1. \] 
    We claim that $|E_s|\le \gamma m$. 
    Suppose on the contrary that $|E_s|> \gamma m$. Then these $\gamma m$ edges destroy at least $\gamma m (|C|/3-1)$ triangles of $T_{C,3}$. So $G$ has at most $(|C|/3 +1)^3 - \gamma m (|C|/3-1)$ good triangles. Since $|C|\le \lceil 3(m/3)^{1/2}\rceil$, it is easy to verify that 
    \[ (|C|/3 +1)^3 - \gamma m (|C|/3-1) < (1- \beta - \gamma^2)(m/3)^{3/2}, 
    \] 
    which leads to a contradiction on the number of good triangles in $G$. 

    Therefore, we conclude that 
    $d(G,T_{C,3})\le |E_b| + |E_s| \le 
    2\gamma m = \varepsilon m$, as desired. 
\end{proof}






\end{document}
